\documentclass{beamer}
\usetheme{Madrid}
\usecolortheme{default}
\usepackage[utf8]{inputenc}
\usepackage{amsmath}
\usepackage{amsfonts}
\usepackage{amssymb}
\usepackage{tikz}

\title{CRISC Exam - Hard Questions with Explanations}
\subtitle{Each Option Analyzed}
\author{Open Source Community }
\date{\today}

\begin{document}
	
	\begin{frame}
		\titlepage
	\end{frame}
	
	\begin{frame}{Outline}
		\tableofcontents
	\end{frame}
	
	\section{Question 1: KRIs Maintenance}
	
	\begin{frame}{Q1: Why Maintain KRIs?}
		\begin{block}{Question}
			Which of the following is the MOST important reason to maintain key risk indicators (KRIs)?
		\end{block}
		\begin{enumerate}[(A)]
			\item In order to avoid risk
			\item Complex metrics require fine-tuning
			\item Risk reports need to be timely
			\item Threats and vulnerabilities change over time
		\end{enumerate}
	\end{frame}
	
	\begin{frame}{Q1: Explanation}
		\begin{block}{Correct Answer: D}
			Threats and vulnerabilities change over time and KRI maintenance ensures that KRIs continue to effectively capture these changes.
		\end{block}
		\begin{block}{Option Analysis}
			\begin{itemize}
				\item \textbf{A - Incorrect:} Risk avoidance is a risk response, not a reason for KRI maintenance
				\item \textbf{B - Incorrect:} Fine-tuning is secondary; capturing changes is primary
				\item \textbf{C - Incorrect:} Timeliness is a business requirement, not a maintenance reason
				\item \textbf{D - Correct:} Risk environment is dynamic; KRIs must evolve with threats
			\end{itemize}
		\end{block}
	\end{frame}
	
	\section{Question 2: Lessons Learned}
	
	\begin{frame}{Q2: Risk Responses After Project}
		\begin{block}{Question}
			What should you do with the risk responses identified during the project's monitoring and controlling process after project completion?
		\end{block}
		\begin{enumerate}[(A)]
			\item Include in project management plan
			\item Include in risk management plan
			\item Include in lessons learned database
			\item Nothing - already in risk register
		\end{enumerate}
	\end{frame}
	
	\begin{frame}{Q2: Explanation}
		\begin{block}{Correct Answer: C}
			Risk responses should be included in the organization's lessons learned database so other project managers can use them.
		\end{block}
		\begin{block}{Option Analysis}
			\begin{itemize}
				\item \textbf{A - Incorrect:} Project management plan is for current project execution
				\item \textbf{B - Incorrect:} Risk management plan is for current project methodology
				\item \textbf{C - Correct:} Knowledge sharing benefits future projects
				\item \textbf{D - Incorrect:} Risk register alone doesn't share knowledge across projects
			\end{itemize}
		\end{block}
	\end{frame}
	
	\section{Question 3: Positive Risk Response}
	
	\begin{frame}{Q3: Risk That Saves Money}
		\begin{block}{Question}
			You identified a risk event that could save \$100,000 in project costs if it occurs. Which statement BEST describes this risk event?
		\end{block}
		\begin{enumerate}[(A)]
			\item Should be mitigated
			\item Should be accepted
			\item Should be avoided
			\item This is an opportunity and should be exploited
		\end{enumerate}
	\end{frame}
	
	\begin{frame}{Q3: Explanation}
		\begin{block}{Correct Answer: D}
			This risk event has potential to save money, so it is an opportunity. The appropriate strategy is to exploit it.
		\end{block}
		\begin{block}{Option Analysis}
			\begin{itemize}
				\item \textbf{A - Incorrect:} Mitigation is for negative risks, not positive opportunities
				\item \textbf{B - Incorrect:} Acceptance is passive; opportunity deserves active exploitation
				\item \textbf{C - Incorrect:} Avoidance is for negative risks, prevents opportunity
				\item \textbf{D - Correct:} Exploiting ensures the opportunity is realized
			\end{itemize}
		\end{block}
	\end{frame}
	
	\section{Question 4: Risk Identification Sessions}
	
	\begin{frame}{Q4: Multiple Risk Identification Meetings}
		\begin{block}{Question}
			Why schedule multiple risk identification meetings throughout the project rather than just initially?
		\end{block}
		\begin{enumerate}[(A)]
			\item Allow all stakeholders to participate throughout phases
			\item Discuss events which did not happen
			\item Identify newly discovered risk events
			\item Communicate pending risks
		\end{enumerate}
	\end{frame}
	
	\begin{frame}{Q4: Explanation}
		\begin{block}{Correct Answer: C}
			Risk identification is iterative because new risks may evolve or become known as the project progresses.
		\end{block}
		\begin{block}{Option Analysis}
			\begin{itemize}
				\item \textbf{A - Incorrect:} Participation is secondary to identifying new risks
				\item \textbf{B - Incorrect:} Focus is on discovering new events, not past events
				\item \textbf{C - Correct:} New risks emerge throughout project life cycle
				\item \textbf{D - Incorrect:} Communication is secondary to identification
			\end{itemize}
		\end{block}
	\end{frame}
	
	\section{Question 5: Risk Priority Number}
	
	\begin{frame}{Q5: RPN Calculation}
		\begin{block}{Question}
			A risk has rating for occurrence, severity, and detection as 4, 5, and 6 respectively. What Risk Priority Number (RPN) would you give?
		\end{block}
		\begin{enumerate}[(A)]
			\item 120
			\item 100
			\item 15
			\item 30
		\end{enumerate}
	\end{frame}
	
	\begin{frame}{Q5: Explanation}
		\begin{block}{Correct Answer: A}
			RPN = Severity × Occurrence × Detection = 4 × 5 × 6 = 120
		\end{block}
		\begin{block}{Option Analysis}
			\begin{itemize}
				\item \textbf{A - Correct:} 4 × 5 × 6 = 120
				\item \textbf{B - Incorrect:} Would be sum of values multiplied incorrectly
				\item \textbf{C - Incorrect:} This would be addition (4+5+6=15)
				\item \textbf{D - Incorrect:} This ignores one factor (4×5=20 or 5×6=30)
			\end{itemize}
		\end{block}
	\end{frame}
	
	\section{Question 6: KRI Most Important Use}
	
	\begin{frame}{Q6: Primary Use of KRIs}
		\begin{block}{Question}
			Which is the MOST important use of KRIs?
		\end{block}
		\begin{enumerate}[(A)]
			\item Backward-looking view of past events
			\item Providing an early warning signal
			\item Indication of risk appetite and tolerance
			\item Enabling documentation of trends
		\end{enumerate}
	\end{frame}
	
	\begin{frame}{Q6: Explanation}
		\begin{block}{Correct Answer: B}
			The most important function of KRIs is to provide an early warning signal that a high risk is emerging.
		\end{block}
		\begin{block}{Option Analysis}
			\begin{itemize}
				\item \textbf{A - Incorrect:} KRIs are forward-looking, not backward-looking
				\item \textbf{B - Correct:} Early warning enables proactive management
				\item \textbf{C - Incorrect:} This is a secondary use, not the most important
				\item \textbf{D - Incorrect:} Documentation is a supporting function
			\end{itemize}
		\end{block}
	\end{frame}
	
	\section{Question 7: KRI Decision Makers}
	
	\begin{frame}{Q7: Who Decides KRIs?}
		\begin{block}{Question}
			Which role carriers decide the Key Risk Indicators of the enterprise? (Choose two)
		\end{block}
		\begin{enumerate}[(A)]
			\item Business leaders
			\item Senior management
			\item Human resources
			\item Chief financial officer
		\end{enumerate}
	\end{frame}
	
	\begin{frame}{Q7: Explanation}
		\begin{block}{Correct Answer: A and B}
			CRISC works with senior management and business leaders to determine which risk indicators become KRIs.
		\end{block}
		\begin{block}{Option Analysis}
			\begin{itemize}
				\item \textbf{A - Correct:} Business leaders provide strategic risk perspective
				\item \textbf{B - Correct:} Senior management approves KRIs
				\item \textbf{C - Incorrect:} HR oversees risk culture, not KRI selection
				\item \textbf{D - Incorrect:} CFO oversees financial risks only
			\end{itemize}
		\end{block}
	\end{frame}
	
	\section{Question 8: Risk Scenario Requirements}
	
	\begin{frame}{Q8: Risk Scenario Requirements}
		\begin{block}{Question}
			What are the requirements for creating risk scenarios? (Choose three)
		\end{block}
		\begin{enumerate}[(A)]
			\item Determination of cause and effect
			\item Determination of value of business process at risk
			\item Potential threats and vulnerabilities that could cause loss
			\item Determination of value of an asset
		\end{enumerate}
	\end{frame}
	
	\begin{frame}{Q8: Explanation}
		\begin{block}{Correct Answer: B, C, D}
			Creating risk scenarios requires: value of asset/business process at risk, and potential threats/vulnerabilities.
		\end{block}
		\begin{block}{Option Analysis}
			\begin{itemize}
				\item \textbf{A - Incorrect:} Cause-effect is predictive analysis, not scenario creation
				\item \textbf{B - Correct:} Business process value is essential for scenario
				\item \textbf{C - Correct:} Threats and vulnerabilities define potential loss
				\item \textbf{D - Correct:} Asset value is fundamental to risk quantification
			\end{itemize}
		\end{block}
	\end{frame}
	
	\section{Question 9: Risk Communications Plan}
	
	\begin{frame}{Q9: Who Shares Risk Information?}
		\begin{block}{Question}
			Which project management plan defines who will be available to share information on project risks?
		\end{block}
		\begin{enumerate}[(A)]
			\item Resource Management Plan
			\item Risk Management Plan
			\item Stakeholder Management Strategy
			\item Communications Management Plan
		\end{enumerate}
	\end{frame}
	
	\begin{frame}{Q9: Explanation}
		\begin{block}{Correct Answer: D}
			The Communications Management Plan defines who will be available to share information on risks and responses throughout the project.
		\end{block}
		\begin{block}{Option Analysis}
			\begin{itemize}
				\item \textbf{A - Incorrect:} Resource Management Plan doesn't define risk communications
				\item \textbf{B - Incorrect:} Risk Management Plan defines analysis and responses, not communication roles
				\item \textbf{C - Incorrect:} Stakeholder strategy doesn't address risk communications
				\item \textbf{D - Correct:} Communications Plan defines roles for sharing risk information
			\end{itemize}
		\end{block}
	\end{frame}
	
	\section{Question 10: Non-Technical Controls}
	
	\begin{frame}{Q10: Non-Technical Controls Example}
		\begin{block}{Question}
			Which control is an example of non-technical controls?
		\end{block}
		\begin{enumerate}[(A)]
			\item Access control
			\item Physical security
			\item Intrusion detection system
			\item Encryption
		\end{enumerate}
	\end{frame}
	
	\begin{frame}{Q10: Explanation}
		\begin{block}{Correct Answer: B}
			Physical security is an example of non-technical control (operational controls family).
		\end{block}
		\begin{block}{Option Analysis}
			\begin{itemize}
				\item \textbf{A - Incorrect:} Access control is technical (system-based)
				\item \textbf{B - Correct:} Physical security is operational, not technical
				\item \textbf{C - Incorrect:} IDS is technical (hardware/software based)
				\item \textbf{D - Incorrect:} Encryption is technical (software/firmware based)
			\end{itemize}
		\end{block}
	\end{frame}
	
	\section{Question 11: Risk Identification Tools}
	
	\begin{frame}{Q11: NOT Used in Risk Identification}
		\begin{block}{Question}
			Which tool is NOT used for identifying project risks?
		\end{block}
		\begin{enumerate}[(A)]
			\item Process flowchart
			\item Ishikawa diagram
			\item Influence diagram
			\item Decision tree diagram
		\end{enumerate}
	\end{frame}
	
	\begin{frame}{Q11: Explanation}
		\begin{block}{Correct Answer: D}
			Decision tree diagrams are used during Quantitative risk analysis, not in risk identification.
		\end{block}
		\begin{block}{Option Analysis}
			\begin{itemize}
				\item \textbf{A - Incorrect:} Process flowchart is used in identification
				\item \textbf{B - Incorrect:} Ishikawa diagram identifies causes
				\item \textbf{C - Incorrect:} Influence diagram shows relationships
				\item \textbf{D - Correct:} Decision tree is for quantitative analysis, not identification
			\end{itemize}
		\end{block}
	\end{frame}
	
	\section{Question 12: Risk Utility}
	
	\begin{frame}{Q12: Definition of Risk Utility}
		\begin{block}{Question}
			Which BEST describes the utility of a risk?
		\end{block}
		\begin{enumerate}[(A)]
			\item The finance incentive behind the risk
			\item The potential opportunity of the risk
			\item The mechanics of how a risk works
			\item The usefulness of the risk to individuals or groups
		\end{enumerate}
	\end{frame}
	
	\begin{frame}{Q12: Explanation}
		\begin{block}{Correct Answer: D}
			Risk utility describes the usefulness of a particular risk to an individual or group.
		\end{block}
		\begin{block}{Option Analysis}
			\begin{itemize}
				\item \textbf{A - Incorrect:} Finance incentive is a measurement method, not definition
				\item \textbf{B - Incorrect:} This describes opportunity, not utility
				\item \textbf{C - Incorrect:} This describes risk mechanics, not utility
				\item \textbf{D - Correct:} Utility = usefulness/benefit to stakeholders
			\end{itemize}
		\end{block}
	\end{frame}
	
	\section{Question 13: Monitoring Tool Scalability}
	
	\begin{frame}{Q13: Monitoring Tool Growth Ability}
		\begin{block}{Question}
			Which monitoring tool aspect ensures ability to keep up with enterprise growth?
		\end{block}
		\begin{enumerate}[(A)]
			\item Scalability
			\item Customizability
			\item Sustainability
			\item Impact on performance
		\end{enumerate}
	\end{frame}
	
	\begin{frame}{Q13: Explanation}
		\begin{block}{Correct Answer: A}
			Scalability ensures monitoring tools can meet anticipated growth in process, complexity, or transaction volumes.
		\end{block}
		\begin{block}{Option Analysis}
			\begin{itemize}
				\item \textbf{A - Correct:} Scalability is about growth capacity
				\item \textbf{B - Incorrect:} Customizability is about adapting to needs
				\item \textbf{C - Incorrect:} Sustainability is about long-term effectiveness
				\item \textbf{D - Incorrect:} Performance impact is separate from growth ability
			\end{itemize}
		\end{block}
	\end{frame}
	
	\section{Question 14: Risk Levels}
	
	\begin{frame}{Q14: Moderate Risk Definition}
		\begin{block}{Question}
			A risk that is a noticeable failure threatening the success of certain goals exists at which level?
		\end{block}
		\begin{enumerate}[(A)]
			\item Moderate risk
			\item High risk
			\item Extremely high risk
			\item Low risk
		\end{enumerate}
	\end{frame}
	
	\begin{frame}{Q14: Explanation}
		\begin{block}{Correct Answer: A}
			Moderate risks are noticeable failures threatening the success of certain goals.
		\end{block}
		\begin{block}{Option Analysis}
			\begin{itemize}
				\item \textbf{A - Correct:} Moderate = noticeable failure affecting specific goals
				\item \textbf{B - Incorrect:} High = significant failure impacting multiple goals
				\item \textbf{C - Incorrect:} Extremely high = large impact, likely failure with severe consequences
				\item \textbf{D - Incorrect:} Low = results in certain unsuccessful goals (minimal)
			\end{itemize}
		\end{block}
	\end{frame}
	
	\section{Question 15: Grouping Risks by Cause}
	
	\begin{frame}{Q15: Grouping Risks by Common Causes}
		\begin{block}{Question}
			What is the primary advantage to group risks by common causes during qualitative risk analysis?
		\end{block}
		\begin{enumerate}[(A)]
			\item Realize areas most laden with risks
			\item Assist in developing effective risk responses
			\item Save time by collecting related resources
			\item Create risk categories unique to each project
		\end{enumerate}
	\end{frame}
	
	\begin{frame}{Q15: Explanation}
		\begin{block}{Correct Answer: B}
			By grouping risks by categories, the project team can develop effective risk responses. Related risks often have common causal factors.
		\end{block}
		\begin{block}{Option Analysis}
			\begin{itemize}
				\item \textbf{A - Incorrect:} This is a byproduct, not primary advantage
				\item \textbf{B - Correct:} Common causes can be addressed with single response
				\item \textbf{C - Incorrect:} Time savings is secondary to response effectiveness
				\item \textbf{D - Incorrect:} Categories may be organizational, not unique
			\end{itemize}
		\end{block}
	\end{frame}
	
	\section{Question 16: Risk Communication Definition}
	
	\begin{frame}{Q16: Risk Communication Process}
		\begin{block}{Question}
			"Process of exchanging information and views about risks among stakeholders" describes which process?
		\end{block}
		\begin{enumerate}[(A)]
			\item Risk governance
			\item Risk identification
			\item Risk response planning
			\item Risk communication
		\end{enumerate}
	\end{frame}
	
	\begin{frame}{Q16: Explanation}
		\begin{block}{Correct Answer: D}
			Risk communication is exactly defined as exchanging information and views about risks among stakeholders.
		\end{block}
		\begin{block}{Option Analysis}
			\begin{itemize}
				\item \textbf{A - Incorrect:} Governance is about decision-making structures
				\item \textbf{B - Incorrect:} Identification is about discovering risks
				\item \textbf{C - Incorrect:} Response planning is about mitigation actions
				\item \textbf{D - Correct:} Communication is information exchange
			\end{itemize}
		\end{block}
	\end{frame}
	
	\section{Question 17: Risk Identification Output}
	
	\begin{frame}{Q17: Key Output of Identify Risk Process}
		\begin{block}{Question}
			Which is a key output of the Identify Risks process?
		\end{block}
		\begin{enumerate}[(A)]
			\item Risk Register
			\item Risk Management Plan
			\item Risk Breakdown Structure
			\item Risk Categories
		\end{enumerate}
	\end{frame}
	
	\begin{frame}{Q17: Explanation}
		\begin{block}{Correct Answer: A}
			The primary output from Identify Risks is the initial entries into the risk register.
		\end{block}
		\begin{block}{Option Analysis}
			\begin{itemize}
				\item \textbf{A - Correct:} Risk register is the only output of Identify Risks
				\item \textbf{B - Incorrect:} Risk Management Plan is from Plan Risk Management
				\item \textbf{C - Incorrect:} RBS is from Plan Risk Management
				\item \textbf{D - Incorrect:} Categories are from Plan Risk Management
			\end{itemize}
		\end{block}
	\end{frame}
	
	\section{Question 18: Risk Scenario Components}
	
	\begin{frame}{Q18: Component Generating Threat}
		\begin{block}{Question}
			Which risk scenario component has potential to generate internal or external threat?
		\end{block}
		\begin{enumerate}[(A)]
			\item Timing dimension
			\item Events
			\item Assets
			\item Actors
		\end{enumerate}
	\end{frame}
	
	\begin{frame}{Q18: Explanation}
		\begin{block}{Correct Answer: D}
			Actors are those components that have the potential to generate the threat (internal or external, human or non-human).
		\end{block}
		\begin{block}{Option Analysis}
			\begin{itemize}
				\item \textbf{A - Incorrect:} Timing dimension is about when, not threat generation
				\item \textbf{B - Incorrect:} Events are happenings, not threat generators
				\item \textbf{C - Incorrect:} Assets are what need protection, not threat sources
				\item \textbf{D - Correct:} Actors generate threats (staff, competitors, outsiders)
			\end{itemize}
		\end{block}
	\end{frame}
	
	\section{Question 19: Control Reliability}
	
	\begin{frame}{Q19: Before Relying on Controls}
		\begin{block}{Question}
			What should you do before relying on any control?
		\end{block}
		\begin{enumerate}[(A)]
			\item Review performance data
			\item Discover risk exposure
			\item Conduct pilot testing
			\item Articulate risk
		\end{enumerate}
	\end{frame}
	
	\begin{frame}{Q19: Explanation}
		\begin{block}{Correct Answer: A and C}
			Pilot testing and reviewing performance data to verify operation against design are done before relying on control.
		\end{block}
		\begin{block}{Option Analysis}
			\begin{itemize}
				\item \textbf{A - Correct:} Performance data verifies control operation
				\item \textbf{B - Incorrect:} Risk exposure helps identify severity, not control reliability
				\item \textbf{C - Correct:} Pilot testing validates control effectiveness
				\item \textbf{D - Incorrect:} Articulating risk is first phase, doesn't verify control
			\end{itemize}
		\end{block}
	\end{frame}
	
	\section{Question 20: Risk Maturity Level 1}
	
	\begin{frame}{Q20: NOT True for Maturity Level 1}
		\begin{block}{Question}
			Which statement is NOT true for risk management capability maturity level 1?
		\end{block}
		\begin{enumerate}[(A)]
			\item Risk is viewed as technical issue
			\item Decisions involving risk lack credible information
			\item Risk appetite applied only during episodic assessments
			\item Risk management skills exist on ad hoc basis
		\end{enumerate}
	\end{frame}
	
	\begin{frame}{Q20: Explanation}
		\begin{block}{Correct Answer: B}
			Decisions lacking credible information is a characteristic of Level 0, not Level 1.
		\end{block}
		\begin{block}{Option Analysis}
			\begin{itemize}
				\item \textbf{A - True:} Level 1 views risk as technical issue
				\item \textbf{B - Not True:} This is Level 0 characteristic
				\item \textbf{C - True:} Level 1 has episodic risk appetite application
				\item \textbf{D - True:} Level 1 has ad hoc risk management skills
			\end{itemize}
		\end{block}
	\end{frame}
	
	\section{Question 21: Reviewing Risk Response Options}
	
	\begin{frame}{Q21: Most Important for Risk Review}
		\begin{block}{Question}
			Which stakeholder is MOST important for reviewing risk response options?
		\end{block}
		\begin{enumerate}[(A)]
			\item Information security managers
			\item Internal auditors
			\item Incident response team members
			\item Business managers
		\end{enumerate}
	\end{frame}
	
	\begin{frame}{Q21: Explanation}
		\begin{block}{Correct Answer: D}
			Business managers are accountable for managing associated risk and determine actions based on information provided by others.
		\end{block}
		\begin{block}{Option Analysis}
			\begin{itemize}
				\item \textbf{A - Incorrect:} Security managers understand technical aspects but business managers decide
				\item \textbf{B - Incorrect:} Auditors provide assurance, not decision-making
				\item \textbf{C - Incorrect:} Incident team responds to events, doesn't review options
				\item \textbf{D - Correct:} Business managers are accountable and decide actions
			\end{itemize}
		\end{block}
	\end{frame}
	
	\section{Question 22: Fault Tree Analysis}
	
	\begin{frame}{Q22: Fault Tree Analysis}
		\begin{block}{Question}
			Which technique provides systematic description of unwanted occurrences in a system?
		\end{block}
		\begin{enumerate}[(A)]
			\item Sensitivity analysis
			\item Scenario analysis
			\item Fault tree analysis
			\item Cause and effect analysis
		\end{enumerate}
	\end{frame}
	
	\begin{frame}{Q22: Explanation}
		\begin{block}{Correct Answer: C}
			Fault Tree Analysis (FTA) provides systematic description of possible occurrences that can result in undesirable outcomes.
		\end{block}
		\begin{block}{Option Analysis}
			\begin{itemize}
				\item \textbf{A - Incorrect:} Sensitivity examines input uncertainty impact
				\item \textbf{B - Incorrect:} Scenario shows range of values across scenarios
				\item \textbf{C - Correct:} FTA combines hardware and human failures
				\item \textbf{D - Incorrect:} Cause-effect explores root causes
			\end{itemize}
		\end{block}
	\end{frame}
	
	\section{Question 23: Risk Treatment}
	
	\begin{frame}{Q23: Risk Treatment Definition}
		\begin{block}{Question}
			What is the process for selecting and implementing measures to impact risk called?
		\end{block}
		\begin{enumerate}[(A)]
			\item Risk Treatment
			\item Control
			\item Risk Assessment
			\item Risk Management
		\end{enumerate}
	\end{frame}
	
	\begin{frame}{Q23: Explanation}
		\begin{block}{Correct Answer: A}
			Risk treatment is the process for selecting and implementing measures for impacting risk in the environment.
		\end{block}
		\begin{block}{Option Analysis}
			\begin{itemize}
				\item \textbf{A - Correct:} Treatment specifically addresses implementation
				\item \textbf{B - Incorrect:} Control is one type of measure, not the process
				\item \textbf{C - Incorrect:} Assessment is analyzing and evaluating, not implementing
				\item \textbf{D - Incorrect:} Management is broader coordinated activities
			\end{itemize}
		\end{block}
	\end{frame}
	
	\section{Question 24: SOX Section 302}
	
	\begin{frame}{Q24: Sarbanes-Oxley Section 302}
		\begin{block}{Question}
			Which SOX section specifies "Periodic financial reports must be certified by CEO and CFO"?
		\end{block}
		\begin{enumerate}[(A)]
			\item Section 302
			\item Section 404
			\item Section 203
			\item Section 409
		\end{enumerate}
	\end{frame}
	
	\begin{frame}{Q24: Explanation}
		\begin{block}{Correct Answer: A}
			Section 302 requires corporate responsibility for financial reports to be certified by CEO, CFO, or designated representative.
		\end{block}
		\begin{block}{Option Analysis}
			\begin{itemize}
				\item \textbf{A - Correct:} Section 302 is certification requirement
				\item \textbf{B - Incorrect:} Section 404 is annual internal controls assessment
				\item \textbf{C - Incorrect:} Section 203 is audit partner rotation
				\item \textbf{D - Incorrect:} Section 409 is timely financial reporting
			\end{itemize}
		\end{block}
	\end{frame}
	
	\section{Question 25: Primary Need for Controls Assessment}
	
	\begin{frame}{Q25: Primary Need for Control Assessment}
		\begin{block}{Question}
			What is the PRIMARY need for effectively assessing controls?
		\end{block}
		\begin{enumerate}[(A)]
			\item Control's alignment with operating environment
			\item Control's design effectiveness
			\item Control's objective achievement
			\item Control's operating effectiveness
		\end{enumerate}
	\end{frame}
	
	\begin{frame}{Q25: Explanation}
		\begin{block}{Correct Answer: C}
			Controls can be effectively assessed only by determining how accurately the control objective is achieved within the operating environment.
		\end{block}
		\begin{block}{Option Analysis}
			\begin{itemize}
				\item \textbf{A - Incorrect:} Alignment is essential but secondary to objective achievement
				\item \textbf{B - Incorrect:} Design effectiveness is considered after objectives
				\item \textbf{C - Correct:} Objective achievement is the top priority
				\item \textbf{D - Incorrect:} Operating effectiveness is considered but after objectives
			\end{itemize}
		\end{block}
	\end{frame}
	
	\section{Question 26: Fast Tracking}
	
	\begin{frame}{Q26: Fast Tracking Effect}
		\begin{block}{Question}
			When you fast track the project, what is likely to increase?
		\end{block}
		\begin{enumerate}[(A)]
			\item Human resource needs
			\item Quality control concerns
			\item Costs
			\item Risks
		\end{enumerate}
	\end{frame}
	
	\begin{frame}{Q26: Explanation}
		\begin{block}{Correct Answer: D}
			Fast tracking allows phases to overlap and generally increases risks within the project.
		\end{block}
		\begin{block}{Option Analysis}
			\begin{itemize}
				\item \textbf{A - Incorrect:} HR is not typically affected by fast tracking
				\item \textbf{B - Incorrect:} Quality concerns usually not affected
				\item \textbf{C - Incorrect:} Costs don't generally increase from fast tracking
				\item \textbf{D - Correct:} Overlapping phases introduces more risk
			\end{itemize}
		\end{block}
	\end{frame}
	
	\section{Question 27: Risk Mitigation}
	
	\begin{frame}{Q27: Identifying Risk Response Type}
		\begin{block}{Question}
			David takes actions to ensure risk event doesn't happen, costing \$10,000. What risk response is this?
		\end{block}
		\begin{enumerate}[(A)]
			\item Avoidance
			\item Mitigation
			\item Acceptance
			\item Transfer
		\end{enumerate}
	\end{frame}
	
	\begin{frame}{Q27: Explanation}
		\begin{block}{Correct Answer: B}
			Taking operational controls to reduce likelihood and impact of risk is risk mitigation.
		\end{block}
		\begin{block}{Option Analysis}
			\begin{itemize}
				\item \textbf{A - Incorrect:} Avoidance discontinues activities causing risk
				\item \textbf{B - Correct:} Actions reduce probability and/or impact
				\item \textbf{C - Incorrect:} Acceptance means no action taken
				\item \textbf{D - Incorrect:} Transfer would involve third party
			\end{itemize}
		\end{block}
	\end{frame}
	
	\section{Question 28: IS Control Objective}
	
	\begin{frame}{Q28: Most Important IS Control Objective}
		\begin{block}{Question}
			What is the MOST important objective of information system control?
		\end{block}
		\begin{enumerate}[(A)]
			\item Business objectives achieved, undesired risk detected/corrected
			\item Ensuring effective and efficient operations
			\item Developing business continuity plans
			\item Safeguarding assets
		\end{enumerate}
	\end{frame}
	
	\begin{frame}{Q28: Explanation}
		\begin{block}{Correct Answer: A}
			The basic purpose of IS control is to ensure business objectives are achieved and undesired risk events are detected and corrected.
		\end{block}
		\begin{block}{Option Analysis}
			\begin{itemize}
				\item \textbf{A - Correct:} Primary purpose - business objectives + risk detection
				\item \textbf{B - Incorrect:} This is one objective but not the most important
				\item \textbf{C - Incorrect:} This is one objective but not the most important
				\item \textbf{D - Incorrect:} This is one objective but not the most important
			\end{itemize}
		\end{block}
	\end{frame}
	
	\section{Question 29: Business Continuity Strategy}
	
	\begin{frame}{Q29: Starting Point for IT Service Continuity}
		\begin{block}{Question}
			Which is prepared by business and serves as starting point for IT Service Continuity Strategy?
		\end{block}
		\begin{enumerate}[(A)]
			\item Business Continuity Strategy
			\item Index of Disaster-Relevant Information
			\item Disaster Invocation Guideline
			\item Testing Schedule
		\end{enumerate}
	\end{frame}
	
	\begin{frame}{Q29: Explanation}
		\begin{block}{Correct Answer: A}
			Business Continuity Strategy is outline of approach to ensure continuity of Vital Business Functions, prepared by business.
		\end{block}
		\begin{block}{Option Analysis}
			\begin{itemize}
				\item \textbf{A - Correct:} Business creates BCS which informs IT strategy
				\item \textbf{B - Incorrect:} Index is IT document for disaster information
				\item \textbf{C - Incorrect:} Guideline is IT document with disaster instructions
				\item \textbf{D - Incorrect:} Testing Schedule is IT operational document
			\end{itemize}
		\end{block}
	\end{frame}
	
	\section{Question 30: Maturity Level 5}
	
	\begin{frame}{Q30: Maturity Level 5 Statement}
		\begin{block}{Question}
			"Real-time monitoring of risk events and control exceptions exists, as does automation of policy management" describes which level?
		\end{block}
		\begin{enumerate}[(A)]
			\item Level 3
			\item Level 0
			\item Level 5
			\item Level 2
		\end{enumerate}
	\end{frame}
	
	\begin{frame}{Q30: Explanation}
		\begin{block}{Correct Answer: C}
			Risk management capability maturity level 5 has real-time monitoring and automation of policy management.
		\end{block}
		\begin{block}{Option Analysis}
			\begin{itemize}
				\item \textbf{A - Incorrect:} Level 3 doesn't have real-time monitoring
				\item \textbf{B - Incorrect:} Level 0 doesn't recognize risk management
				\item \textbf{C - Correct:} Level 5 has real-time monitoring and automation
				\item \textbf{D - Incorrect:} Level 2 doesn't have real-time monitoring
			\end{itemize}
		\end{block}
	\end{frame}
	
	\section{Question 31: Cost Performance Index}
	
	\begin{frame}{Q31: CPI True Statement}
		\begin{block}{Question}
			Which is true for Cost Performance Index (CPI)?
		\end{block}
		\begin{enumerate}[(A)]
			\item If CPI > 1, it indicates better than expected performance
			\item CPI = EV × AC
			\item CPI is used to measure schedule performance
			\item If CPI = 1, it indicates poor performance
		\end{enumerate}
	\end{frame}
	
	\begin{frame}{Q31: Explanation}
		\begin{block}{Correct Answer: A}
			CPI = EV/AC. If CPI > 1, indicates better than expected performance. CPI = 1 indicates right on target.
		\end{block}
		\begin{block}{Option Analysis}
			\begin{itemize}
				\item \textbf{A - Correct:} >1 = better than expected performance
				\item \textbf{B - Incorrect:} CPI = EV/AC, not EV × AC
				\item \textbf{C - Incorrect:} CPI measures cost performance, not schedule
				\item \textbf{D - Incorrect:} CPI = 1 indicates right on target
			\end{itemize}
		\end{block}
	\end{frame}
	
	\section{Question 32: Direct vs Indirect Information}
	
	\begin{frame}{Q32: NOT Indirect Information}
		\begin{block}{Question}
			Which does NOT provide indirect information?
		\end{block}
		\begin{enumerate}[(A)]
			\item Information about propriety of cutoff
			\item Reports of orders rejected for credit limitations
			\item Reports of unusual deviations
			\item No significant differences between perpetual and actual levels
		\end{enumerate}
	\end{frame}
	
	\begin{frame}{Q32: Explanation}
		\begin{block}{Correct Answer: A}
			Information about propriety of cutoff is direct information, not indirect.
		\end{block}
		\begin{block}{Option Analysis}
			\begin{itemize}
				\item \textbf{A - Correct:} Propriety of cutoff is direct evidence
				\item \textbf{B - Incorrect:} Rejected orders gives indirect info about credit controls
				\item \textbf{C - Incorrect:} Unusual deviations gives indirect info about billing controls
				\item \textbf{D - Incorrect:} No differences gives indirect info about billing controls
			\end{itemize}
		\end{block}
	\end{frame}
	
	\section{Question 33: Salience Model}
	
	\begin{frame}{Q33: Salience Model Description}
		\begin{block}{Question}
			Which activity best describes a salience model?
		\end{block}
		\begin{enumerate}[(A)]
			\item Power, urgency, and legitimacy
			\item Power and interest
			\item Influence and impact
			\item Power and influence
		\end{enumerate}
	\end{frame}
	
	\begin{frame}{Q33: Explanation}
		\begin{block}{Correct Answer: A}
			Salience model describes stakeholders based on power (ability to impose will), urgency (need for attention), and legitimacy (appropriate involvement).
		\end{block}
		\begin{block}{Option Analysis}
			\begin{itemize}
				\item \textbf{A - Correct:} Power, urgency, legitimacy = salience model
				\item \textbf{B - Incorrect:} This is power/interest grid
				\item \textbf{C - Incorrect:} This is influence/impact grid
				\item \textbf{D - Incorrect:} This is power/influence grid
			\end{itemize}
		\end{block}
	\end{frame}
	
	\section{Question 34: First Step in Risk Assessment}
	
	\begin{frame}{Q34: First Step in Risk Assessment}
		\begin{block}{Question}
			What is the first step in the risk assessment process?
		\end{block}
		\begin{enumerate}[(A)]
			\item Identification of assets
			\item Identification of threats
			\item Identification of threat sources
			\item Identification of vulnerabilities
		\end{enumerate}
	\end{frame}
	
	\begin{frame}{Q34: Explanation}
		\begin{block}{Correct Answer: A}
			Asset identification is the most crucial and first step in the risk assessment process.
		\end{block}
		\begin{block}{Option Analysis}
			\begin{itemize}
				\item \textbf{A - Correct:} Assets come first (people, processes, infrastructure)
				\item \textbf{B - Incorrect:} Threats come after asset identification
				\item \textbf{C - Incorrect:} Threat sources come after assets
				\item \textbf{D - Incorrect:} Vulnerabilities come after threat identification
			\end{itemize}
		\end{block}
	\end{frame}
	
	\section{Question 35: Risk Threshold Matrix}
	
	\begin{frame}{Q35: Risk Threshold Specification}
		\begin{block}{Question}
			Which matrix is used to specify risk thresholds?
		\end{block}
		\begin{enumerate}[(A)]
			\item Risk indicator matrix
			\item Impact matrix
			\item Risk scenario matrix
			\item Probability matrix
		\end{enumerate}
	\end{frame}
	
	\begin{frame}{Q35: Explanation}
		\begin{block}{Correct Answer: A}
			Risk indicators are metrics used to indicate risk thresholds - when risk level is approaching high or unacceptable level.
		\end{block}
		\begin{block}{Option Analysis}
			\begin{itemize}
				\item \textbf{A - Correct:} Risk indicators specifically show thresholds
				\item \textbf{B - Incorrect:} Impact matrix helps rating risks, not thresholds
				\item \textbf{C - Incorrect:} Scenario matrix describes events, not thresholds
				\item \textbf{D - Incorrect:} Probability matrix helps rating risks, not thresholds
			\end{itemize}
		\end{block}
	\end{frame}
	
	\section{Question 36: Risk Appetite Factors}
	
	\begin{frame}{Q36: Major Risk Appetite Factors}
		\begin{block}{Question}
			What are the two MAJOR factors for deciding risk appetite? (Choose two)
		\end{block}
		\begin{enumerate}[(A)]
			\item Amount of loss enterprise wants to accept
			\item Alignment with risk-culture
			\item Risk-aware decisions
			\item Capacity to absorb loss
		\end{enumerate}
	\end{frame}
	
	\begin{frame}{Q36: Explanation}
		\begin{block}{Correct Answer: A and D}
			Two major factors: (1) capacity to absorb loss, (2) culture towards risk taking - amount of loss enterprise wants to accept.
		\end{block}
		\begin{block}{Option Analysis}
			\begin{itemize}
				\item \textbf{A - Correct:} Amount of loss to accept is key factor
				\item \textbf{B - Incorrect:} Alignment with risk-culture is secondary
				\item \textbf{C - Incorrect:} Risk-aware decisions are results, not factors
				\item \textbf{D - Correct:} Capacity to absorb loss is key factor
			\end{itemize}
		\end{block}
	\end{frame}
	
	\section{Question 37: Risk Monitoring Process}
	
	\begin{frame}{Q37: Process Responsible for Risk Updates}
		\begin{block}{Question}
			Which process is responsible for identifying new risks, changing risks, and outdated risks?
		\end{block}
		\begin{enumerate}[(A)]
			\item Risk planning
			\item Risk monitoring and controlling
			\item Risk identification
			\item Risk analysis
		\end{enumerate}
	\end{frame}
	
	\begin{frame}{Q37: Explanation}
		\begin{block}{Correct Answer: B}
			Risk monitoring and controlling is responsible for identifying new risks, determining status of changed risks, and outdated risks.
		\end{block}
		\begin{block}{Option Analysis}
			\begin{itemize}
				\item \textbf{A - Incorrect:} Planning creates methodology
				\item \textbf{B - Correct:} Monitoring tracks risk status changes
				\item \textbf{C - Incorrect:} Identification happens before analysis
				\item \textbf{D - Incorrect:} Analysis determines severity, not dynamic changes
			\end{itemize}
		\end{block}
	\end{frame}
	
	\section{Question 38: Single Loss Expectancy}
	
	\begin{frame}{Q38: SLE Calculation}
		\begin{block}{Question}
			Asset valued at \$125,000 subjected to exposure factor of 25\%. What is Single Loss Expectancy?
		\end{block}
		\begin{enumerate}[(A)]
			\item \$125,025
			\item \$31,250
			\item \$5,000
			\item \$3,125,000
		\end{enumerate}
	\end{frame}
	
	\begin{frame}{Q38: Explanation}
		\begin{block}{Correct Answer: B}
			SLE = Asset Value × Exposure Factor = \$125,000 × 0.25 = \$31,250
		\end{block}
		\begin{block}{Option Analysis}
			\begin{itemize}
				\item \textbf{A - Incorrect:} \$125,025 is adding instead of multiplying
				\item \textbf{B - Correct:} \$125,000 × 0.25 = \$31,250
				\item \textbf{C - Incorrect:} \$5,000 = incorrect calculation
				\item \textbf{D - Incorrect:} \$3,125,000 is dividing instead of multiplying
			\end{itemize}
		\end{block}
	\end{frame}
	
	\section{Question 39: Access Control Principles}
	
	\begin{frame}{Q39: Access Control Principles}
		\begin{block}{Question}
			Which are principles of access controls? (Choose three)
		\end{block}
		\begin{enumerate}[(A)]
			\item Confidentiality
			\item Availability
			\item Reliability
			\item Integrity
		\end{enumerate}
	\end{frame}
	
	\begin{frame}{Q39: Explanation}
		\begin{block}{Correct Answer: A, B, D}
			Access controls focus on Confidentiality, Integrity, and Availability (CIA triad).
		\end{block}
		\begin{block}{Option Analysis}
			\begin{itemize}
				\item \textbf{A - Correct:} Confidentiality is core principle
				\item \textbf{B - Correct:} Availability is core principle
				\item \textbf{C - Incorrect:} Reliability is not a CIA principle
				\item \textbf{D - Correct:} Integrity is core principle
			\end{itemize}
		\end{block}
	\end{frame}
	
	\section{Question 40: KRI Maintenance Reason}
	
	\begin{frame}{Q40: Maintaining KRIs Reason}
		\begin{block}{Question}
			What is the MOST important reason to maintain Key Risk Indicators?
		\end{block}
		\begin{enumerate}[(A)]
			\item Risk reports need to be timely
			\item Complex metrics require fine-tuning
			\item Threats and vulnerabilities change over time
			\item They help to avoid risk
		\end{enumerate}
	\end{frame}
	
	\begin{frame}{Q40: Explanation}
		\begin{block}{Correct Answer: C}
			Since internal and external environments constantly change, KRIs need to effectively capture these changes.
		\end{block}
		\begin{block}{Option Analysis}
			\begin{itemize}
				\item \textbf{A - Incorrect:} Timeliness is business requirement, not maintenance reason
				\item \textbf{B - Incorrect:} Fine-tuning is secondary to capturing changes
				\item \textbf{C - Correct:} Threats and vulnerabilities change over time
				\item \textbf{D - Incorrect:} Avoiding risk is a response, not maintenance reason
			\end{itemize}
		\end{block}
	\end{frame}
	
	\section{Question 41: Technical Controls}
	
	\begin{frame}{Q41: NOT Technical Control}
		\begin{block}{Question}
			Which control does NOT come under technical class of control?
		\end{block}
		\begin{enumerate}[(A)]
			\item Program management control
			\item System and Communications Protection control
			\item Identification and Authentication control
			\item Access Control
		\end{enumerate}
	\end{frame}
	
	\begin{frame}{Q41: Explanation}
		\begin{block}{Correct Answer: A}
			Program Management control comes under management class of controls, not technical.
		\end{block}
		\begin{block}{Option Analysis}
			\begin{itemize}
				\item \textbf{A - Correct:} Program Management is management class
				\item \textbf{B - Incorrect:} SC family is technical class
				\item \textbf{C - Incorrect:} IA family is technical class
				\item \textbf{D - Incorrect:} AC family is technical class
			\end{itemize}
		\end{block}
	\end{frame}
	
	\section{Question 42: Brainstorming}
	
	\begin{frame}{Q42: Risk Identification Method}
		\begin{block}{Question}
			Mary uses a facilitator to help generate ideas about project risks. What method is she using?
		\end{block}
		\begin{enumerate}[(A)]
			\item Delphi Techniques
			\item Expert judgment
			\item Brainstorming
			\item Checklist analysis
		\end{enumerate}
	\end{frame}
	
	\begin{frame}{Q42: Explanation}
		\begin{block}{Correct Answer: C}
			Brainstorming attempts to create a comprehensive list of risks and is often led by a moderator or facilitator.
		\end{block}
		\begin{block}{Option Analysis}
			\begin{itemize}
				\item \textbf{A - Incorrect:} Delphi uses anonymous surveys, no facilitator
				\item \textbf{B - Incorrect:} Expert judgment is individual assessment
				\item \textbf{C - Correct:} Brainstorming uses facilitator for idea generation
				\item \textbf{D - Incorrect:} Checklist uses historical information
			\end{itemize}
		\end{block}
	\end{frame}
	
	\section{Question 43: Administrative Control}
	
	\begin{frame}{Q43: Administrative Control Example}
		\begin{block}{Question}
			Which is an administrative control?
		\end{block}
		\begin{enumerate}[(A)]
			\item Water detection
			\item Reasonableness check
			\item Data loss prevention program
			\item Session timeout
		\end{enumerate}
	\end{frame}
	
	\begin{frame}{Q43: Explanation}
		\begin{block}{Correct Answer: C}
			Data Loss Prevention program is an administrative control (program consists of projects, policies).
		\end{block}
		\begin{block}{Option Analysis}
			\begin{itemize}
				\item \textbf{A - Incorrect:} Water detection is physical/technical
				\item \textbf{B - Incorrect:} Reasonableness check is technical
				\item \textbf{C - Correct:} DLP program is administrative
				\item \textbf{D - Incorrect:} Session timeout is technical
			\end{itemize}
		\end{block}
	\end{frame}
	
	\section{Question 44: Risk Management Plan}
	
	\begin{frame}{Q44: Risk Management Plan Creation}
		\begin{block}{Question}
			What document defines how risks will be identified, quantified, and contingency plans implemented?
		\end{block}
		\begin{enumerate}[(A)]
			\item Project plan
			\item Resource management plan
			\item Project management plan
			\item Risk management plan
		\end{enumerate}
	\end{frame}
	
	\begin{frame}{Q44: Explanation}
		\begin{block}{Correct Answer: D}
			Risk management plan defines how risks will be identified, analyzed, monitored, controlled, and responded to.
		\end{block}
		\begin{block}{Option Analysis}
			\begin{itemize}
				\item \textbf{A - Incorrect:} Project plan is not official PMBOK plan
				\item \textbf{B - Incorrect:} Resource plan manages people/equipment
				\item \textbf{C - Incorrect:} PM plan is overarching, not risk-specific
				\item \textbf{D - Correct:} Risk management plan specifically addresses risk
			\end{itemize}
		\end{block}
	\end{frame}
	
	\section{Question 45: Risk Documentation}
	
	\begin{frame}{Q45: Where Risks and Responses are Documented}
		\begin{block}{Question}
			Where are all risks and risk responses documented as project progresses?
		\end{block}
		\begin{enumerate}[(A)]
			\item Risk management plan
			\item Project management plan
			\item Risk response plan
			\item Risk register
		\end{enumerate}
	\end{frame}
	
	\begin{frame}{Q45: Explanation}
		\begin{block}{Correct Answer: D}
			All risks, their responses, and other characteristics are documented in the risk register.
		\end{block}
		\begin{block}{Option Analysis}
			\begin{itemize}
				\item \textbf{A - Incorrect:} Risk Management Plan is methodology, not documentation
				\item \textbf{B - Incorrect:} PM plan is overarching
				\item \textbf{C - Incorrect:} Risk response plan addresses planned responses only
				\item \textbf{D - Correct:} Risk register is the repository for all risk information
			\end{itemize}
		\end{block}
	\end{frame}
	
	\section{Question 46: Risk Transference}
	
	\begin{frame}{Q46: Hiring Third Party}
		\begin{block}{Question}
			Hiring a company to deal with all hardware work is which type of risk response?
		\end{block}
		\begin{enumerate}[(A)]
			\item Transference
			\item Mitigation
			\item Avoidance
			\item Exploit
		\end{enumerate}
	\end{frame}
	
	\begin{frame}{Q46: Explanation}
		\begin{block}{Correct Answer: A}
			Hiring a third party to own risk is risk transference.
		\end{block}
		\begin{block}{Option Analysis}
			\begin{itemize}
				\item \textbf{A - Correct:} Transfer risk to another party
				\item \textbf{B - Incorrect:} Mitigation is spending money to reduce risk
				\item \textbf{C - Incorrect:} Avoidance adds activities to prevent risk
				\item \textbf{D - Incorrect:} Exploit is for positive risks/opportunities
			\end{itemize}
		\end{block}
	\end{frame}
	
	\section{Question 47: Activity Duration Estimates}
	
	\begin{frame}{Q47: Risk Identification Input}
		\begin{block}{Question}
			Which input of identify risks process helps identify risks associated with time allowances?
		\end{block}
		\begin{enumerate}[(A)]
			\item Activity duration estimates
			\item Activity cost estimates
			\item Risk management plan
			\item Schedule management plan
		\end{enumerate}
	\end{frame}
	
	\begin{frame}{Q47: Explanation}
		\begin{block}{Correct Answer: A}
			Activity duration estimates review helps identify risks associated with time allowances, with width of range indicating risk degree.
		\end{block}
		\begin{block}{Option Analysis}
			\begin{itemize}
				\item \textbf{A - Correct:} Duration estimates identify time risks
				\item \textbf{B - Incorrect:} Cost estimates identify cost risks
				\item \textbf{C - Incorrect:} Risk plan identifies methodology
				\item \textbf{D - Incorrect:} Schedule plan describes reporting
			\end{itemize}
		\end{block}
	\end{frame}
	
	\section{Question 48: Loss of Integrity}
	
	\begin{frame}{Q48: Loss of Integrity Events}
		\begin{block}{Question}
			Which events refer to loss of integrity? (Choose three)
		\end{block}
		\begin{enumerate}[(A)]
			\item Someone sees company's secret formula
			\item Someone makes unauthorized changes to a Web site
			\item An e-mail message is modified in transit
			\item A virus infects a file
		\end{enumerate}
	\end{frame}
	
	\begin{frame}{Q48: Explanation}
		\begin{block}{Correct Answer: B, C, D}
			Loss of integrity includes: message modified in transit, virus infects file, unauthorized Web site changes.
		\end{block}
		\begin{block}{Option Analysis}
			\begin{itemize}
				\item \textbf{A - Incorrect:} Seeing secret formula = loss of confidentiality
				\item \textbf{B - Correct:} Unauthorized changes = integrity loss
				\item \textbf{C - Correct:} Message modified = integrity loss
				\item \textbf{D - Correct:} Virus infects file = integrity loss
			\end{itemize}
		\end{block}
	\end{frame}
	
	\section{Question 49: IS Control Design}
	
	\begin{frame}{Q49: First Step in IS Control Design}
		\begin{block}{Question}
			What should be PRIMARILY considered while designing information systems controls?
		\end{block}
		\begin{enumerate}[(A)]
			\item The IT strategic plan
			\item The existing IT environment
			\item The organizational strategic plan
			\item The present IT budget
		\end{enumerate}
	\end{frame}
	
	\begin{frame}{Q49: Explanation}
		\begin{block}{Correct Answer: C}
			Review of enterprise's strategic plan is the first step in designing effective IS controls that fit enterprise's long-term plans.
		\end{block}
		\begin{block}{Option Analysis}
			\begin{itemize}
				\item \textbf{A - Incorrect:} IT plan supports enterprise plan, not primary
				\item \textbf{B - Incorrect:} Existing environment review is useful but secondary
				\item \textbf{C - Correct:} Organizational strategic plan is primary
				\item \textbf{D - Incorrect:} Budget is one component, not primary
			\end{itemize}
		\end{block}
	\end{frame}
	
	\section{Question 50: Effective Communication Inhibitor}
	
	\begin{frame}{Q50: Most Effective Communication Inhibitor}
		\begin{block}{Question}
			Which is the MOST effective inhibitor of relevant and efficient communication?
		\end{block}
		\begin{enumerate}[(A)]
			\item False sense of confidence at top
			\item Perception of covering up known risk
			\item Existence of a blame culture
			\item Misalignment between risk appetite and policies
		\end{enumerate}
	\end{frame}
	
	\begin{frame}{Q50: Explanation}
		\begin{block}{Correct Answer: C}
			Blame culture is the most effective inhibitor. Business units point fingers at IT when projects fail.
		\end{block}
		\begin{block}{Option Analysis}
			\begin{itemize}
				\item \textbf{A - Incorrect:} This is consequence of poor communication
				\item \textbf{B - Incorrect:} This is consequence of poor communication
				\item \textbf{C - Correct:} Blame culture directly inhibits collaboration
				\item \textbf{D - Incorrect:} Misalignment is inhibitor but not as prominent
			\end{itemize}
		\end{block}
	\end{frame}
	
	\section{Question 51: Low Priority Risks}
	
	\begin{frame}{Q51: Handling Small Risks}
		\begin{block}{Question}
			What should you do with small risks that won't affect the project much?
		\end{block}
		\begin{enumerate}[(A)]
			\item These risks can be dismissed
			\item These risks can be accepted
			\item Added to low priority risk watch list
			\item All risks must have valid documented risk response
		\end{enumerate}
	\end{frame}
	
	\begin{frame}{Q51: Explanation}
		\begin{block}{Correct Answer: C}
			Low-impact, low-probability risks can be added to the low priority risk watch list.
		\end{block}
		\begin{block}{Option Analysis}
			\begin{itemize}
				\item \textbf{A - Incorrect:} Risks are not dismissed; they are documented
				\item \textbf{B - Incorrect:} They may be accepted but should be documented on watch list
				\item \textbf{C - Correct:} Watch list for future monitoring
				\item \textbf{D - Incorrect:} Not every risk demands a response
			\end{itemize}
		\end{block}
	\end{frame}
	
	\section{Question 52: IDS Control Type}
	
	\begin{frame}{Q52: Intrusion Detection System Type}
		\begin{block}{Question}
			What type of control is an intrusion detection system (IDS)?
		\end{block}
		\begin{enumerate}[(A)]
			\item Detective
			\item Corrective
			\item Preventative
			\item Recovery
		\end{enumerate}
	\end{frame}
	
	\begin{frame}{Q52: Explanation}
		\begin{block}{Correct Answer: A}
			IDS detects and gives warning when policy violation occurs; it is a detective control.
		\end{block}
		\begin{block}{Option Analysis}
			\begin{itemize}
				\item \textbf{A - Correct:} IDS detects and alerts
				\item \textbf{B - Incorrect:} Corrective reduces impact after detection
				\item \textbf{C - Incorrect:} Preventive would stop before occurrence
				\item \textbf{D - Incorrect:} Recovery overcomes impact on business
			\end{itemize}
		\end{block}
	\end{frame}
	
	\section{Question 53: Audit and Accountability}
	
	\begin{frame}{Q53: Audit and Accountability Functions}
		\begin{block}{Question}
			What are functions of audit and accountability control? (Choose three)
		\end{block}
		\begin{enumerate}[(A)]
			\item Provides details on how to protect audit logs
			\item Implement effective access control
			\item Implement an effective audit program
			\item Provides details on how to determine what to audit
		\end{enumerate}
	\end{frame}
	
	\begin{frame}{Q53: Explanation}
		\begin{block}{Correct Answer: A, C, D}
			Audit and accountability helps: implement audit program, determine what to audit, protect audit logs, non-repudiation.
		\end{block}
		\begin{block}{Option Analysis}
			\begin{itemize}
				\item \textbf{A - Correct:} Protects audit logs
				\item \textbf{B - Incorrect:} Access control is separate family
				\item \textbf{C - Correct:} Implements audit program
				\item \textbf{D - Correct:} Determines what to audit
			\end{itemize}
		\end{block}
	\end{frame}
	
	\section{Question 54: Risk Response Trigger}
	
	\begin{frame}{Q54: Trigger for Risk Response}
		\begin{block}{Question}
			Which acts as a trigger for risk response process?
		\end{block}
		\begin{enumerate}[(A)]
			\item Risk level increases above risk appetite
			\item Risk level increases above risk tolerance
			\item Risk level equates risk appetite
			\item Risk level equates risk tolerance
		\end{enumerate}
	\end{frame}
	
	\begin{frame}{Q54: Explanation}
		\begin{block}{Correct Answer: B}
			Risk response process is triggered when a risk exceeds the enterprise's risk tolerance level.
		\end{block}
		\begin{block}{Option Analysis}
			\begin{itemize}
				\item \textbf{A - Incorrect:} Appetite is not the trigger
				\item \textbf{B - Correct:} Tolerance exceeded = trigger
				\item \textbf{C - Incorrect:} Equating appetite is not trigger
				\item \textbf{D - Incorrect:} Equating tolerance is not trigger
			\end{itemize}
		\end{block}
	\end{frame}
	
	\section{Question 55: Exposure Factor Value}
	
	\begin{frame}{Q55: Complete Asset Loss}
		\begin{block}{Question}
			What is the value of exposure factor if the asset is lost completely?
		\end{block}
		\begin{enumerate}[(A)]
			\item 1
			\item Infinity
			\item 10
			\item 0
		\end{enumerate}
	\end{frame}
	
	\begin{frame}{Q55: Explanation}
		\begin{block}{Correct Answer: A}
			Exposure Factor = 1.0 when asset is completely lost (100\% loss).
		\end{block}
		\begin{block}{Option Analysis}
			\begin{itemize}
				\item \textbf{A - Correct:} 1.0 = 100\% loss
				\item \textbf{B - Incorrect:} Infinity is not used
				\item \textbf{C - Incorrect:} 10 is not valid
				\item \textbf{D - Incorrect:} 0 = 0\% loss
			\end{itemize}
		\end{block}
	\end{frame}
	
	\section{Question 56: Positive Risk Exploit}
	
	\begin{frame}{Q56: Exploiting Positive Risk}
		\begin{block}{Question}
			Discovering a profitable byproduct from your project is an example of what risk response?
		\end{block}
		\begin{enumerate}[(A)]
			\item Enhancing
			\item Positive
			\item Opportunistic
			\item Exploiting
		\end{enumerate}
	\end{frame}
	
	\begin{frame}{Q56: Explanation}
		\begin{block}{Correct Answer: D}
			This is exploiting a positive risk - a byproduct is an excellent example of exploiting an opportunity.
		\end{block}
		\begin{block}{Option Analysis}
			\begin{itemize}
				\item \textbf{A - Incorrect:} Enhancing increases probability, not seizing
				\item \textbf{B - Incorrect:} Positive describes risk type, not response
				\item \textbf{C - Incorrect:} Opportunistic is not a valid risk response
				\item \textbf{D - Correct:} Exploiting seizes the opportunity
			\end{itemize}
		\end{block}
	\end{frame}
	
	\section{Question 57: ALE Formula}
	
	\begin{frame}{Q57: ALE Formula}
		\begin{block}{Question}
			Which is true for SLE, ARO, and ALE?
		\end{block}
		\begin{enumerate}[(A)]
			\item ALE = ARO / SLE
			\item ARO = SLE / ALE
			\item ARO = ALE × SLE
			\item ALE = ARO × SLE
		\end{enumerate}
	\end{frame}
	
	\begin{frame}{Q57: Explanation}
		\begin{block}{Correct Answer: D}
			ALE = SLE × ARO
		\end{block}
		\begin{block}{Option Analysis}
			\begin{itemize}
				\item \textbf{A - Incorrect:} Wrong formula
				\item \textbf{B - Incorrect:} Wrong formula
				\item \textbf{C - Incorrect:} Wrong formula
				\item \textbf{D - Correct:} ALE = ARO × SLE
			\end{itemize}
		\end{block}
	\end{frame}
	
	\section{Question 58: Maturity Level 3}
	
	\begin{frame}{Q58: Maturity Level 3 Statements}
		\begin{block}{Question}
			Which statements are true for risk management capability maturity level 3?
		\end{block}
		\begin{enumerate}[(A)]
			\item Workflow tools used to accelerate risk issues
			\item Business knows IT fits in enterprise risk universe
			\item Formal continuous improvement required
			\item Risk viewed as business issue
		\end{enumerate}
	\end{frame}
	
	\begin{frame}{Q58: Explanation}
		\begin{block}{Correct Answer: A, B, D}
			Level 3: risk viewed as business issue, business knows IT fits in risk universe, workflow tools used.
		\end{block}
		\begin{block}{Option Analysis}
			\begin{itemize}
				\item \textbf{A - True:} Workflow tools used in Level 3
				\item \textbf{B - True:} IT fits in risk universe at Level 3
				\item \textbf{C - Not True:} This is Level 5
				\item \textbf{D - True:} Risk viewed as business issue at Level 3
			\end{itemize}
		\end{block}
	\end{frame}
	
	\section{Question 59: Risk Profile Responsibility}
	
	\begin{frame}{Q59: Analyzing Risks Responsibility}
		\begin{block}{Question}
			Which role carrier is accounted for analyzing risks, maintaining risk profile, and risk-aware decisions?
		\end{block}
		\begin{enumerate}[(A)]
			\item Business management
			\item Business process owner
			\item Chief information officer
			\item Chief risk officer
		\end{enumerate}
	\end{frame}
	
	\begin{frame}{Q59: Explanation}
		\begin{block}{Correct Answer: A}
			Business management is accountable for analyzing risks, maintaining risk profile, and risk-aware decisions.
		\end{block}
		\begin{block}{Option Analysis}
			\begin{itemize}
				\item \textbf{A - Correct:} Business management accountable
				\item \textbf{B - Incorrect:} Process owner responsible but not accountable
				\item \textbf{C - Incorrect:} CIO has some responsibility but not accountability
				\item \textbf{D - Incorrect:} CRO has oversight, not direct accountability
			\end{itemize}
		\end{block}
	\end{frame}
	
	\section{Question 60: Poor Password Quality}
	
	\begin{frame}{Q60: Poor Password and Unsafe Transmission}
		\begin{block}{Question}
			Choosing poor password and transmitting data unprotected refers to?
		\end{block}
		\begin{enumerate}[(A)]
			\item Probabilities
			\item Threats
			\item Vulnerabilities
			\item Impacts
		\end{enumerate}
	\end{frame}
	
	\begin{frame}{Q60: Explanation}
		\begin{block}{Correct Answer: C}
			Vulnerabilities are characteristics that may be exploited by a threat. Poor password and unprotected transmission are vulnerabilities.
		\end{block}
		\begin{block}{Option Analysis}
			\begin{itemize}
				\item \textbf{A - Incorrect:} Probability is likelihood of occurrence
				\item \textbf{B - Incorrect:} Threat is potential cause of harm
				\item \textbf{C - Correct:} Vulnerabilities are weaknesses
				\item \textbf{D - Incorrect:} Impact is result of exploitation
			\end{itemize}
		\end{block}
	\end{frame}
	
	\section{Question 61: Outsourced Service Provider Compliance}
	
	\begin{frame}{Q61: Ensuring Outsourced Compliance}
		\begin{block}{Question}
			Which is BEST way to ensure outsourced service providers comply with security policy?
		\end{block}
		\begin{enumerate}[(A)]
			\item Penetration testing
			\item Service level monitoring
			\item Security awareness training
			\item Periodic audits
		\end{enumerate}
	\end{frame}
	
	\begin{frame}{Q61: Explanation}
		\begin{block}{Correct Answer: D}
			Periodic audits can spot gaps in information security compliance.
		\end{block}
		\begin{block}{Option Analysis}
			\begin{itemize}
				\item \textbf{A - Incorrect:} Penetration testing identifies vulnerability, not compliance
				\item \textbf{B - Incorrect:} Service monitoring identifies operational issues
				\item \textbf{C - Incorrect:} Training increases awareness but less effective
				\item \textbf{D - Correct:} Periodic audits ensure compliance
			\end{itemize}
		\end{block}
	\end{frame}
	
	\section{Question 62: Risk Prioritization}
	
	\begin{frame}{Q62: Risk Prioritization Category}
		\begin{block}{Question}
			Re-architecting existing system due to complexity - which risk prioritization option?
		\end{block}
		\begin{enumerate}[(A)]
			\item Deferrals
			\item Quick win
			\item Business case to be made
			\item Contagious risk
		\end{enumerate}
	\end{frame}
	
	\begin{frame}{Q62: Explanation}
		\begin{block}{Correct Answer: C}
			This is categorized as business case to be made because the response requires quite large investment.
		\end{block}
		\begin{block}{Option Analysis}
			\begin{itemize}
				\item \textbf{A - Incorrect:} Deferrals address costly response to low risk
				\item \textbf{B - Incorrect:} Quick win requires small investment
				\item \textbf{C - Correct:} Large investment = business case to be made
				\item \textbf{D - Incorrect:} Contagious risk is a type of risk, not prioritization
			\end{itemize}
		\end{block}
	\end{frame}
	
	\section{Question 63: Firewall Compliance}
	
	\begin{frame}{Q63: Ensuring Firewall Compliance}
		\begin{block}{Question}
			What BEST ensures firewall configured in compliance with security policy?
		\end{block}
		\begin{enumerate}[(A)]
			\item Interview the firewall administrator
			\item Review the actual procedures
			\item Review device's log file
			\item Review the parameter settings
		\end{enumerate}
	\end{frame}
	
	\begin{frame}{Q63: Explanation}
		\begin{block}{Correct Answer: D}
			Review of parameter settings provides basis for comparison of actual configuration to security policy.
		\end{block}
		\begin{block}{Option Analysis}
			\begin{itemize}
				\item \textbf{A - Incorrect:} Interview provides overview, not confirmation
				\item \textbf{B - Incorrect:} Procedures show how it should be managed, not compliance
				\item \textbf{C - Incorrect:} Log review shows logging is enabled, not configuration
				\item \textbf{D - Correct:} Parameter settings show actual configuration
			\end{itemize}
		\end{block}
	\end{frame}
	
	\section{Question 64: Critical Success Factors}
	
	\begin{frame}{Q64: NOT Used for CSF Measurement}
		\begin{block}{Question}
			Which is NOT used for measurement of Critical Success Factors?
		\end{block}
		\begin{enumerate}[(A)]
			\item Productivity
			\item Quality
			\item Quantity
			\item Customer service
		\end{enumerate}
	\end{frame}
	
	\begin{frame}{Q64: Explanation}
		\begin{block}{Correct Answer: C}
			Quantity is NOT used for measuring Critical Success Factors.
		\end{block}
		\begin{block}{Option Analysis}
			\begin{itemize}
				\item \textbf{A - Incorrect:} Productivity is used
				\item \textbf{B - Incorrect:} Quality is used
				\item \textbf{C - Correct:} Quantity is not used
				\item \textbf{D - Incorrect:} Customer service is used
			\end{itemize}
		\end{block}
	\end{frame}
	
	\section{Question 65: Risk Management Plan}
	
	\begin{frame}{Q65: Risk Management Plan Statement}
		\begin{block}{Question}
			Which statement is NOT true regarding risk management plan?
		\end{block}
		\begin{enumerate}[(A)]
			\item Output of Plan Risk Management
			\item Input to all remaining risk-planning
			\item Includes description of risk responses and triggers
			\item Includes thresholds, scoring, responsible parties, budgets
		\end{enumerate}
	\end{frame}
	
	\begin{frame}{Q65: Explanation}
		\begin{block}{Correct Answer: C}
			Risk management plan does NOT include responses to risks or triggers. Responses are in risk register.
		\end{block}
		\begin{block}{Option Analysis}
			\begin{itemize}
				\item \textbf{A - True:} Output of Plan Risk Management
				\item \textbf{B - True:} Input to other processes
				\item \textbf{C - Not True:} Responses/triggers are in risk register
				\item \textbf{D - True:} Includes thresholds, scoring, parties, budgets
			\end{itemize}
		\end{block}
	\end{frame}
	
	\section{Question 66: Risk Response Tool}
	
	\begin{frame}{Q66: Choosing Risk Response Tool}
		\begin{block}{Question}
			Which tool would you use to choose appropriate risk response?
		\end{block}
		\begin{enumerate}[(A)]
			\item Project network diagrams
			\item Cause-and-effect analysis
			\item Decision tree analysis
			\item Delphi Technique
		\end{enumerate}
	\end{frame}
	
	\begin{frame}{Q66: Explanation}
		\begin{block}{Correct Answer: C}
			Decision tree analysis helps determine best risk response by measuring probability, impact, and risk exposure.
		\end{block}
		\begin{block}{Option Analysis}
			\begin{itemize}
				\item \textbf{A - Incorrect:} Network diagrams visualize workflow
				\item \textbf{B - Incorrect:} Cause-effect exposes risk factors
				\item \textbf{C - Correct:} Decision tree evaluates response options
				\item \textbf{D - Incorrect:} Delphi identifies most probable risks
			\end{itemize}
		\end{block}
	\end{frame}
	
	\section{Question 67: Maturity Level 0}
	
	\begin{frame}{Q67: Risk Management Maturity Level}
		\begin{block}{Question}
			Enterprise takes decisions without risk credential information - which maturity level?
		\end{block}
		\begin{enumerate}[(A)]
			\item Level 1
			\item Level 0
			\item Level 5
			\item Level 4
		\end{enumerate}
	\end{frame}
	
	\begin{frame}{Q67: Explanation}
		\begin{block}{Correct Answer: B}
			Level 0: enterprise does not recognize need to consider risk management, decisions lack credible information.
		\end{block}
		\begin{block}{Option Analysis}
			\begin{itemize}
				\item \textbf{A - Incorrect:} Level 1 uses credible information
				\item \textbf{B - Correct:} Level 0 lacks risk awareness and information
				\item \textbf{C - Incorrect:} Level 5 has full automation
				\item \textbf{D - Incorrect:} Level 4 integrates with ERM
			\end{itemize}
		\end{block}
	\end{frame}
	
	\section{Question 68: Data Owner Priority}
	
	\begin{frame}{Q68: Data Owner Priority}
		\begin{block}{Question}
			Which is priority of data owners when establishing risk mitigation method?
		\end{block}
		\begin{enumerate}[(A)]
			\item User entitlement changes
			\item Platform security
			\item Intrusion detection
			\item Antivirus controls
		\end{enumerate}
	\end{frame}
	
	\begin{frame}{Q68: Explanation}
		\begin{block}{Correct Answer: A}
			Data owners are responsible for assigning user entitlement changes and approving access to systems.
		\end{block}
		\begin{block}{Option Analysis}
			\begin{itemize}
				\item \textbf{A - Correct:} User entitlement changes = data owner responsibility
				\item \textbf{B - Incorrect:} Platform security = data custodian
				\item \textbf{C - Incorrect:} Intrusion detection = data custodian
				\item \textbf{D - Incorrect:} Antivirus controls = data custodian
			\end{itemize}
		\end{block}
	\end{frame}
	
	\section{Question 69: Acceptable Use Policy}
	
	\begin{frame}{Q69: Forbidding Personal Email Use}
		\begin{block}{Question}
			What policy forbids employees from using organizational e-mail for personal use?
		\end{block}
		\begin{enumerate}[(A)]
			\item Anti-harassment policy
			\item Acceptable use policy
			\item Intellectual property policy
			\item Privacy policy
		\end{enumerate}
	\end{frame}
	
	\begin{frame}{Q69: Explanation}
		\begin{block}{Correct Answer: B}
			Acceptable Use Policy (AUP) restricts ways network, website or system may be used.
		\end{block}
		\begin{block}{Option Analysis}
			\begin{itemize}
				\item \textbf{A - Incorrect:} Anti-harassment is not IS security
				\item \textbf{B - Correct:} AUP defines acceptable system usage
				\item \textbf{C - Incorrect:} Intellectual property is not personal email
				\item \textbf{D - Incorrect:} Privacy policy is about data handling
			\end{itemize}
		\end{block}
	\end{frame}
	
	\section{Question 70: Risk Response Type}
	
	\begin{frame}{Q70: Risk Response Identification}
		\begin{block}{Question}
			Wendy reduces risk impact from \$75,000 to \$15,000 with 10\% chance, costing \$25,000. What risk response?
		\end{block}
		\begin{enumerate}[(A)]
			\item Mitigation
			\item Avoidance
			\item Transference
			\item Enhancing
		\end{enumerate}
	\end{frame}
	
	\begin{frame}{Q70: Explanation}
		\begin{block}{Correct Answer: A}
			Mitigation implies reduction in probability and/or impact of adverse risk event within acceptable thresholds.
		\end{block}
		\begin{block}{Option Analysis}
			\begin{itemize}
				\item \textbf{A - Correct:} Reduction of impact and probability
				\item \textbf{B - Incorrect:} Avoidance changes plan to avoid entirely
				\item \textbf{C - Incorrect:} Transference shifts to third party
				\item \textbf{D - Incorrect:} Enhancing is for positive risks
			\end{itemize}
		\end{block}
	\end{frame}
	
	\section{Question 71: Plan Risk Response}
	
	\begin{frame}{Q71: Plan Risk Response Process}
		\begin{block}{Question}
			Which process addresses risks by priorities, schedules PM plan as required, inserts resources into budget?
		\end{block}
		\begin{enumerate}[(A)]
			\item Monitor and Control Risk
			\item Plan risk response
			\item Identify Risks
			\item Qualitative Risk Analysis
		\end{enumerate}
	\end{frame}
	
	\begin{frame}{Q71: Explanation}
		\begin{block}{Correct Answer: B}
			Plan Risk Response addresses risks by priorities, schedules PM plan, and inserts resources and activities into budget.
		\end{block}
		\begin{block}{Option Analysis}
			\begin{itemize}
				\item \textbf{A - Incorrect:} Monitor and Control implements responses
				\item \textbf{B - Correct:} Plan Risk Response prioritizes and schedules
				\item \textbf{C - Incorrect:} Identify Risks determines risks
				\item \textbf{D - Incorrect:} Qualitative Analysis assesses nature of risks
			\end{itemize}
		\end{block}
	\end{frame}
	
	\section{Question 72: Negative Risk Response}
	
	\begin{frame}{Q72: Negative Risk Response}
		\begin{block}{Question}
			Which risk response is used for negative risk events?
		\end{block}
		\begin{enumerate}[(A)]
			\item Share
			\item Enhance
			\item Exploit
			\item Accept
		\end{enumerate}
	\end{frame}
	
	\begin{frame}{Q72: Explanation}
		\begin{block}{Correct Answer: D}
			Acceptance is used for negative risk events. Share, Enhance, Exploit are for positive risks.
		\end{block}
		\begin{block}{Option Analysis}
			\begin{itemize}
				\item \textbf{A - Incorrect:} Share is for positive risks
				\item \textbf{B - Incorrect:} Enhance is for positive risks
				\item \textbf{C - Incorrect:} Exploit is for positive risks
				\item \textbf{D - Correct:} Accept is for negative risks
			\end{itemize}
		\end{block}
	\end{frame}
	
	\section{Question 73: Investment/Expense Risk}
	
	\begin{frame}{Q73: Investment Risk Definition}
		\begin{block}{Question}
			Which risk refers to probability that actual return is lower than investor's expectations?
		\end{block}
		\begin{enumerate}[(A)]
			\item Integrity risk
			\item Project ownership risk
			\item Relevance risk
			\item Expense risk
		\end{enumerate}
	\end{frame}
	
	\begin{frame}{Q73: Explanation}
		\begin{block}{Correct Answer: D}
			Expense risk is probability that actual return on investment will be lower than investor's expectations.
		\end{block}
		\begin{block}{Option Analysis}
			\begin{itemize}
				\item \textbf{A - Incorrect:} Integrity = data cannot be relied on
				\item \textbf{B - Incorrect:} Project ownership = lack of accountability
				\item \textbf{C - Incorrect:} Relevance = wrong info to wrong people
				\item \textbf{D - Correct:} Expense risk = lower than expected return
			\end{itemize}
		\end{block}
	\end{frame}
	
	\section{Question 74: Risk Scenario Requirements}
	
	\begin{frame}{Q74: PRIMARY Requirements for Risk Scenarios}
		\begin{block}{Question}
			What are PRIMARY requirements for developing risk scenarios? (Choose two)
		\end{block}
		\begin{enumerate}[(A)]
			\item Potential threats and vulnerabilities
			\item Determination of value of asset at risk
			\item Determination of actors
			\item Determination of threat type
		\end{enumerate}
	\end{frame}
	
	\begin{frame}{Q74: Explanation}
		\begin{block}{Correct Answer: A and B}
			Primary requirements: value of asset/business process at risk, and potential threats and vulnerabilities.
		\end{block}
		\begin{block}{Option Analysis}
			\begin{itemize}
				\item \textbf{A - Correct:} Threats and vulnerabilities are primary
				\item \textbf{B - Correct:} Asset value is primary
				\item \textbf{C - Incorrect:} Actors are secondary components
				\item \textbf{D - Incorrect:} Threat type is secondary component
			\end{itemize}
		\end{block}
	\end{frame}
	
	\section{Question 75: CRO Responsibilities}
	
	\begin{frame}{Q75: CRO Responsibilities}
		\begin{block}{Question}
			What are responsibilities of the CRO? (Choose three)
		\end{block}
		\begin{enumerate}[(A)]
			\item Managing risk assessment process
			\item Implement corrective actions
			\item Advising Board of Directors
			\item Managing supporting risk management function
		\end{enumerate}
	\end{frame}
	
	\begin{frame}{Q75: Explanation}
		\begin{block}{Correct Answer: A, B, D}
			CRO responsibilities: managing risk assessment, implementing corrective actions, supporting risk management functions.
		\end{block}
		\begin{block}{Option Analysis}
			\begin{itemize}
				\item \textbf{A - Correct:} Managing risk assessment process
				\item \textbf{B - Correct:} Implementation of corrective actions
				\item \textbf{C - Incorrect:} CRO advises but this is not primary responsibility listed
				\item \textbf{D - Correct:} Supporting risk management functions
			\end{itemize}
		\end{block}
	\end{frame}
	
	\section{Question 76: Vendor Change Control}
	
	\begin{frame}{Q76: Vendor Change Control System}
		\begin{block}{Question}
			Which system introduces and executes stakeholder change request with vendor?
		\end{block}
		\begin{enumerate}[(A)]
			\item Contract change control system
			\item Scope change control system
			\item Cost change control system\item Schedule change control system
		\end{enumerate}
	\end{frame}
	
	\begin{frame}{Q76: Explanation}
		\begin{block}{Correct Answer: A}
			Contract change control system addresses changes with vendor that may affect the project contract.
		\end{block}
		\begin{block}{Option Analysis}
			\begin{itemize}
				\item \textbf{A - Correct:} Contract change control for vendor changes
				\item \textbf{B - Incorrect:} Scope may change but vendor relationship is contract
				\item \textbf{C - Incorrect:} Cost change system manages cost changes
				\item \textbf{D - Incorrect:} No schedule impact indicated
			\end{itemize}
		\end{block}
	\end{frame}
	
	\section{Question 77: Control Cost vs Benefit}
	
	\begin{frame}{Q77: When Control Costs Exceed Benefits}
		\begin{block}{Question}
			Costs of specific controls exceed benefits of mitigating a given risk. What is BEST action?
		\end{block}
		\begin{enumerate}[(A)]
			\item Apply the appropriate control anyway
			\item Adopt corrective control
			\item Accept the risk rather than incur cost of mitigation
			\item Exploit the risk
		\end{enumerate}
	\end{frame}
	
	\begin{frame}{Q77: Explanation}
		\begin{block}{Correct Answer: C}
			If control costs exceed benefits, enterprise may accept risk rather than incur cost of mitigation.
		\end{block}
		\begin{block}{Option Analysis}
			\begin{itemize}
				\item \textbf{A - Incorrect:} Controls are not applied when cost exceeds benefit
				\item \textbf{B - Incorrect:} No control should be adopted
				\item \textbf{C - Correct:} Accept risk when mitigation cost exceeds benefit
				\item \textbf{D - Incorrect:} Exploit is for opportunities, not negative risk
			\end{itemize}
		\end{block}
	\end{frame}
	
	\section{Question 78: Mortality Tables}
	
	\begin{frame}{Q78: Mortality Tables Mathematics}
		\begin{block}{Question}
			Mortality tables are based on what mathematical activity? (Choose three)
		\end{block}
		\begin{enumerate}[(A)]
			\item Normal distributions
			\item Probabilities
			\item Impact
			\item Sampling
		\end{enumerate}
	\end{frame}
	
	\begin{frame}{Q78: Explanation}
		\begin{block}{Correct Answer: A, B, D}
			Mortality tables use: normal distributions, probabilities, and sampling. Impact is not mathematical.
		\end{block}
		\begin{block}{Option Analysis}
			\begin{itemize}
				\item \textbf{A - Correct:} Normal distributions used
				\item \textbf{B - Correct:} Probabilities used
				\item \textbf{C - Incorrect:} Impact uses words (Low/Medium/High), not mathematical
				\item \textbf{D - Correct:} Sampling used
			\end{itemize}
		\end{block}
	\end{frame}
	
	\section{Question 79: Risk Transference}
	
	\begin{frame}{Q79: Hiring Professional Vendor}
		\begin{block}{Question}
			Harry hires professional vendor to complete dangerous portion of work. What risk response?
		\end{block}
		\begin{enumerate}[(A)]
			\item Transference
			\item Mitigation
			\item Acceptance
			\item Avoidance
		\end{enumerate}
	\end{frame}
	
	\begin{frame}{Q79: Explanation}
		\begin{block}{Correct Answer: A}
			Risk transfer shifts impact to another entity. Risk doesn't go away but responsibility is transferred.
		\end{block}
		\begin{block}{Option Analysis}
			\begin{itemize}
				\item \textbf{A - Correct:} Transfer to professional vendor
				\item \textbf{B - Incorrect:} Mitigation would reduce risk in-house
				\item \textbf{C - Incorrect:} Acceptance would retain risk
				\item \textbf{D - Incorrect:} Avoidance removes risk entirely
			\end{itemize}
		\end{block}
	\end{frame}
	
	\section{Question 80: Project Team Involvement}
	
	\begin{frame}{Q80: Why Project Team in Risk Identification}
		\begin{block}{Question}
			Why should project team members be involved in Identify Risk process?
		\end{block}
		\begin{enumerate}[(A)]
			\item They will most likely cause risk events
			\item They will have best responses
			\item They are most affected
			\item They will need sense of ownership and responsibility
		\end{enumerate}
	\end{frame}
	
	\begin{frame}{Q80: Explanation}
		\begin{block}{Correct Answer: D}
			Project team develops sense of ownership and responsibility for risk events and associated responses.
		\end{block}
		\begin{block}{Option Analysis}
			\begin{itemize}
				\item \textbf{A - Incorrect:} Not the primary reason
				\item \textbf{B - Incorrect:} Not the primary reason
				\item \textbf{C - Incorrect:} Not the primary reason
				\item \textbf{D - Correct:} Ownership and responsibility
			\end{itemize}
		\end{block}
	\end{frame}
	
	\section{Question 81: Risk Monitoring Requirements}
	
	\begin{frame}{Q81: Risk Monitoring Requirements}
		\begin{block}{Question}
			What are requirements of monitoring risk? (Choose three)
		\end{block}
		\begin{enumerate}[(A)]
			\item Information of various stakeholders
			\item Preparation of detailed monitoring plan
			\item Identifying the risk to be monitored
			\item Defining the project's scope
		\end{enumerate}
	\end{frame}
	
	\begin{frame}{Q81: Explanation}
		\begin{block}{Correct Answer: B, C, D}
			Requirements: identify risk to monitor, prepare detailed plan, define project scope.
		\end{block}
		\begin{block}{Option Analysis}
			\begin{itemize}
				\item \textbf{A - Incorrect:} Stakeholder info not required for monitoring
				\item \textbf{B - Correct:} Detailed monitoring plan
				\item \textbf{C - Correct:} Identify risk to monitor
				\item \textbf{D - Correct:} Define project scope
			\end{itemize}
		\end{block}
	\end{frame}
	
	\section{Question 82: Insurance Risk Technique}
	
	\begin{frame}{Q82: Insurance Risk Management Technique}
		\begin{block}{Question}
			Liability insurance policy for security risks uses which risk management technique?
		\end{block}
		\begin{enumerate}[(A)]
			\item Risk transfer
			\item Risk acceptance
			\item Risk avoidance
			\item Risk mitigation
		\end{enumerate}
	\end{frame}
	
	\begin{frame}{Q82: Explanation}
		\begin{block}{Correct Answer: A}
			Insurance transfers security risks from company to insurance company.
		\end{block}
		\begin{block}{Option Analysis}
			\begin{itemize}
				\item \textbf{A - Correct:} Transfer risk to insurance
				\item \textbf{B - Incorrect:} Acceptance retains risk
				\item \textbf{C - Incorrect:} Avoidance not performing activity
				\item \textbf{D - Incorrect:} Mitigation reduces severity/likelihood
			\end{itemize}
		\end{block}
	\end{frame}
	
	\section{Question 83: Improving Project Performance}
	
	\begin{frame}{Q83: Improving Performance Through Risk Analysis}
		\begin{block}{Question}
			How can you improve project performance through risk analysis with stakeholders?
		\end{block}
		\begin{enumerate}[(A)]
			\item Involve subject matter experts
			\item Involve stakeholders only in relevant phases
			\item Use qualitative risk analysis quickly
			\item Focus on high-priority risks through qualitative analysis
		\end{enumerate}
	\end{frame}
	
	\begin{frame}{Q83: Explanation}
		\begin{block}{Correct Answer: D}
			By focusing on high-priority risk events through qualitative risk analysis, you can improve project performance.
		\end{block}
		\begin{block}{Option Analysis}
			\begin{itemize}
				\item \textbf{A - Incorrect:} SMEs help but not best approach
				\item \textbf{B - Incorrect:} Stakeholders should be involved throughout
				\item \textbf{C - Incorrect:} Qualitative is fast but not best answer
				\item \textbf{D - Correct:} Focus on high-priority risks
			\end{itemize}
		\end{block}
	\end{frame}
	
	\section{Question 84: Project Risk Timing}
	
	\begin{frame}{Q84: When Risk Events Occur}
		\begin{block}{Question}
			Which statement best addresses when a project risk actually happens?
		\end{block}
		\begin{enumerate}[(A)]
			\item Risks are uncertain as to when they will happen
			\item Risks can happen at any time
			\item Project risks are always in the future
			\item Risk triggers are warning signs
		\end{enumerate}
	\end{frame}
	
	\begin{frame}{Q84: Explanation}
		\begin{block}{Correct Answer: C}
			According to PMBOK, a project risk is always in the future. If risk has happened, it's an issue.
		\end{block}
		\begin{block}{Option Analysis}
			\begin{itemize}
				\item \textbf{A - Incorrect:} Risks are uncertain but can be identified
				\item \textbf{B - Incorrect:} Risks only happen in the future
				\item \textbf{C - Correct:} Risks are always in the future
				\item \textbf{D - Incorrect:} This is true but not best answer
			\end{itemize}
		\end{block}
	\end{frame}
	
	\section{Question 85: Risk Level Indication}
	
	\begin{frame}{Q85: Indication of High Risk Level}
		\begin{block}{Question}
			Which is MOST effective method for indicating risk level approaching unacceptable level?
		\end{block}
		\begin{enumerate}[(A)]
			\item Risk register
			\item Cause and effect diagram
			\item Risk indicator
			\item Return on investment
		\end{enumerate}
	\end{frame}
	
	\begin{frame}{Q85: Explanation}
		\begin{block}{Correct Answer: C}
			Risk indicators are metrics used to indicate risk thresholds - when risk level is approaching high level.
		\end{block}
		\begin{block}{Option Analysis}
			\begin{itemize}
				\item \textbf{A - Incorrect:} Register is inventory, not indicator
				\item \textbf{B - Incorrect:} Cause-effect identifies root causes
				\item \textbf{C - Correct:} Risk indicators show thresholds
				\item \textbf{D - Incorrect:} ROI is performance measure
			\end{itemize}
		\end{block}
	\end{frame}
	
	\section{Question 86: Communications Management Plan}
	
	\begin{frame}{Q86: Risk Information Sharing}
		\begin{block}{Question}
			Which plan defines who will be available to share information on project risks?
		\end{block}
		\begin{enumerate}[(A)]
			\item Risk Management Plan
			\item Stakeholder management strategy
			\item Communications Management Plan
			\item Resource Management Plan
		\end{enumerate}
	\end{frame}
	
	\begin{frame}{Q86: Explanation}
		\begin{block}{Correct Answer: C}
			Communications Management Plan defines who will be available to share information on risks and responses.
		\end{block}
		\begin{block}{Option Analysis}
			\begin{itemize}
				\item \textbf{A - Incorrect:} Risk plan defines identification/analysis
				\item \textbf{B - Incorrect:} Stakeholder strategy doesn't address
				\item \textbf{C - Correct:} Communications Plan defines roles
				\item \textbf{D - Incorrect:} Resource plan doesn't define risk communications
			\end{itemize}
		\end{block}
	\end{frame}
	
	\section{Question 87: Delphi Technique}
	
	\begin{frame}{Q87: Anonymous Risk Identification}
		\begin{block}{Question}
			Which method allows participants to anonymously identify risk events?
		\end{block}
		\begin{enumerate}[(A)]
			\item Delphi technique
			\item Root cause analysis
			\item Isolated pilot groups
			\item SWOT analysis
		\end{enumerate}
	\end{frame}
	
	\begin{frame}{Q87: Explanation}
		\begin{block}{Correct Answer: A}
			Delphi technique uses rounds of anonymous surveys to build consensus and prevents individuals from influencing others.
		\end{block}
		\begin{block}{Option Analysis}
			\begin{itemize}
				\item \textbf{A - Correct:} Anonymous Delphi technique
				\item \textbf{B - Incorrect:} Root cause is not anonymous
				\item \textbf{C - Incorrect:} Isolated pilot groups is not valid
				\item \textbf{D - Incorrect:} SWOT is not anonymous
			\end{itemize}
		\end{block}
	\end{frame}
	
	\section{Question 88: Lack of Controls}
	
	\begin{frame}{Q88: Lack of Adequate Controls}
		\begin{block}{Question}
			Which represents lack of adequate controls?
		\end{block}
		\begin{enumerate}[(A)]
			\item Vulnerability
			\item Threat
			\item Asset
			\item Impact
		\end{enumerate}
	\end{frame}
	
	\begin{frame}{Q88: Explanation}
		\begin{block}{Correct Answer: A}
			Vulnerability is a weakness or lack of safeguard that can be exploited by a threat.
		\end{block}
		\begin{block}{Option Analysis}
			\begin{itemize}
				\item \textbf{A - Correct:} Vulnerability = lack of adequate controls
				\item \textbf{B - Incorrect:} Threat is potential cause
				\item \textbf{C - Incorrect:} Asset is economic resource
				\item \textbf{D - Incorrect:} Impact is financial loss measure
			\end{itemize}
		\end{block}
	\end{frame}
	
	\section{Question 89: Risk Register Updates}
	
	\begin{frame}{Q89: NOT in Risk Register Update}
		\begin{block}{Question}
			Which is NOT included in risk register updates from qualitative risk analysis?
		\end{block}
		\begin{enumerate}[(A)]
			\item Trends in qualitative risk analysis
			\item Risk probability-impact matrix
			\item Risks grouped by categories
			\item Watchlist of low-priority risks
		\end{enumerate}
	\end{frame}
	
	\begin{frame}{Q89: Explanation}
		\begin{block}{Correct Answer: B}
			Risk probability-impact matrix is NOT included in risk register updates. The matrix is a tool.
		\end{block}
		\begin{block}{Option Analysis}
			\begin{itemize}
				\item \textbf{A - Incorrect:} Trends are included
				\item \textbf{B - Correct:} Matrix is not included
				\item \textbf{C - Incorrect:} Categories are included
				\item \textbf{D - Incorrect:} Watchlist is included
			\end{itemize}
		\end{block}
	\end{frame}
	
	\section{Question 90: Systemic Risk}
	
	\begin{frame}{Q90: Systemic Risk}
		\begin{block}{Question}
			Risk that happens with important business partner and affects large group of enterprises?
		\end{block}
		\begin{enumerate}[(A)]
			\item Contagious risk
			\item Reporting risk
			\item Operational risk
			\item Systemic risk
		\end{enumerate}
	\end{frame}
	
	\begin{frame}{Q90: Explanation}
		\begin{block}{Correct Answer: D}
			Systemic risks happen with important business partner and affect large group of enterprises within an area or industry.
		\end{block}
		\begin{block}{Option Analysis}
			\begin{itemize}
				\item \textbf{A - Incorrect:} Contagious = risk events with several partners
				\item \textbf{B - Incorrect:} Reporting = bad decisions from wrong reporting
				\item \textbf{C - Incorrect:} Operational = day-to-day operations risk
				\item \textbf{D - Correct:} Systemic = large group affected
			\end{itemize}
		\end{block}
	\end{frame}
	
	\section{Question 91: Risk Rating Rules}
	
	\begin{frame}{Q91: Risk Prioritization Tool}
		\begin{block}{Question}
			What do you need next to prioritize risks after calculating probability and impact?
		\end{block}
		\begin{enumerate}[(A)]
			\item Affinity Diagram
			\item Risk rating rules
			\item Project Network Diagram
			\item Risk categories
		\end{enumerate}
	\end{frame}
	
	\begin{frame}{Q91: Explanation}
		\begin{block}{Correct Answer: B}
			Risk rating rules define how to prioritize risks after probability and impact values are calculated.
		\end{block}
		\begin{block}{Option Analysis}
			\begin{itemize}
				\item \textbf{A - Incorrect:} Affinity sorts ideas
				\item \textbf{B - Correct:} Risk rating rules prioritize
				\item \textbf{C - Incorrect:} Network shows sequencing
				\item \textbf{D - Incorrect:} Categories are output, not tool
			\end{itemize}
		\end{block}
	\end{frame}
	
	\section{Question 92: Change Request Immediate Action}
	
	\begin{frame}{Q92: Change Request Immediate Action}
		\begin{block}{Question}
			Customer requests change in existing project plan. What is your immediate action?
		\end{block}
		\begin{enumerate}[(A)]
			\item Update risk register
			\item Ask for formal change request
			\item Ignore request
			\item Refuse change request
		\end{enumerate}
	\end{frame}
	
	\begin{frame}{Q92: Explanation}
		\begin{block}{Correct Answer: B}
			Whenever customer requests change, ask for formal change request.
		\end{block}
		\begin{block}{Option Analysis}
			\begin{itemize}
				\item \textbf{A - Incorrect:} First action is formal change request
				\item \textbf{B - Correct:} Require formal change request
				\item \textbf{C - Incorrect:} Cannot ignore
				\item \textbf{D - Incorrect:} Cannot refuse without analysis
			\end{itemize}
		\end{block}
	\end{frame}
	
	\section{Question 93: Certainty Equivalent Value}
	
	\begin{frame}{Q93: Certainty Equivalent Value}
		\begin{block}{Question}
			"Expected guaranteed value of taking a risk" describes what?
		\end{block}
		\begin{enumerate}[(A)]
			\item Certainty equivalent value
			\item Risk premium
			\item Risk value guarantee
			\item Certain value assurance
		\end{enumerate}
	\end{frame}
	
	\begin{frame}{Q93: Explanation}
		\begin{block}{Correct Answer: A}
			Certainty equivalent value is the expected guaranteed value of taking a risk.
		\end{block}
		\begin{block}{Option Analysis}
			\begin{itemize}
				\item \textbf{A - Correct:} Certainty equivalent value
				\item \textbf{B - Incorrect:} Risk premium is difference between expected value and certainty equivalent
				\item \textbf{C - Incorrect:} Not a valid term
				\item \textbf{D - Incorrect:} Not a valid term
			\end{itemize}
		\end{block}
	\end{frame}
	
	\section{Question 94: Risk Trigger}
	
	\begin{frame}{Q94: Risk Trigger}
		\begin{block}{Question}
			Hardware vendor calls saying delivery won't arrive on time. What is this?
		\end{block}
		\begin{enumerate}[(A)]
			\item Residual risk
			\item Trigger
			\item Contingency plan
			\item Secondary risk
		\end{enumerate}
	\end{frame}
	
	\begin{frame}{Q94: Explanation}
		\begin{block}{Correct Answer: B}
			Triggers are warning signs of upcoming risk event. Delay in delivery is a trigger.
		\end{block}
		\begin{block}{Option Analysis}
			\begin{itemize}
				\item \textbf{A - Incorrect:} Residual remains after applying controls
				\item \textbf{B - Correct:} Trigger is warning sign
				\item \textbf{C - Incorrect:} No contingency plan mentioned
				\item \textbf{D - Incorrect:} Secondary from implementing response
			\end{itemize}
		\end{block}
	\end{frame}
	
	\section{Question 95: Quantitative Analysis Inputs}
	
	\begin{frame}{Q95: NOT an Input to Quantitative Analysis}
		\begin{block}{Question}
			Which is NOT an input to quantitative risk analysis process?
		\end{block}
		\begin{enumerate}[(A)]
			\item Risk management plan
			\item Enterprise environmental factors
			\item Cost management plan
			\item Risk register
		\end{enumerate}
	\end{frame}
	
	\begin{frame}{Q95: Explanation}
		\begin{block}{Correct Answer: B}
			Enterprise environmental factors is NOT an input to quantitative risk analysis.
		\end{block}
		\begin{block}{Option Analysis}
			\begin{itemize}
				\item \textbf{A - Incorrect:} Risk management plan is input
				\item \textbf{B - Correct:} EEF is not an input
				\item \textbf{C - Incorrect:} Cost management plan is input
				\item \textbf{D - Incorrect:} Risk register is input
			\end{itemize}
		\end{block}
	\end{frame}
	
	\section{Question 96: Low Risk Events}
	
	\begin{frame}{Q96: Low Probability and Impact Risks}
		\begin{block}{Question}
			What should Stephen do with 47 risks that have low probability and low impact?
		\end{block}
		\begin{enumerate}[(A)]
			\item Accept the risks
			\item Add to watchlist for future monitoring
			\item Dismiss the risks
			\item Add to risk register
		\end{enumerate}
	\end{frame}
	
	\begin{frame}{Q96: Explanation}
		\begin{block}{Correct Answer: B}
			Low probability and low impact risks should be added to a watchlist for future monitoring.
		\end{block}
		\begin{block}{Option Analysis}
			\begin{itemize}
				\item \textbf{A - Incorrect:} They may be accepted but best is watchlist
				\item \textbf{B - Correct:} Watchlist for future monitoring
				\item \textbf{C - Incorrect:} Risks are not dismissed
				\item \textbf{D - Incorrect:} Best is watchlist first
			\end{itemize}
		\end{block}
	\end{frame}
	
	\section{Question 97: Newly Identified Risks}
	
	\begin{frame}{Q97: Newly Identified Risks}
		\begin{block}{Question}
			Jenny uncovers several risk events during quantitative risk analysis. What should she do?
		\end{block}
		\begin{enumerate}[(A)]
			\item Enter into qualitative analysis
			\item Determine if accepted or responded to
			\item Enter into risk register
			\item Continue with quantitative analysis
		\end{enumerate}
	\end{frame}
	
	\begin{frame}{Q97: Explanation}
		\begin{block}{Correct Answer: C}
			All identified risk events should be entered into the risk register.
		\end{block}
		\begin{block}{Option Analysis}
			\begin{itemize}
				\item \textbf{A - Incorrect:} Document first
				\item \textbf{B - Incorrect:} Document and analyze
				\item \textbf{C - Correct:} Enter in risk register
				\item \textbf{D - Incorrect:} Need identification first
			\end{itemize}
		\end{block}
	\end{frame}
	
	\section{Question 98: Risk Avoidance}
	
	\begin{frame}{Q98: Not Engaging in Activity}
		\begin{block}{Question}
			Enterprise chooses not to engage in e-commerce. What risk form is this?
		\end{block}
		\begin{enumerate}[(A)]
			\item Risk avoidance
			\item Risk treatment
			\item Risk acceptance
			\item Risk transfer
		\end{enumerate}
	\end{frame}
	
	\begin{frame}{Q98: Explanation}
		\begin{block}{Correct Answer: A}
			Not engaging in any activity avoids the inherent risk associated with the activity.
		\end{block}
		\begin{block}{Option Analysis}
			\begin{itemize}
				\item \textbf{A - Correct:} Avoidance = not performing activity
				\item \textbf{B - Incorrect:} Treatment = reducing frequency/impact
				\item \textbf{C - Incorrect:} Acceptance = informed decision to accept
				\item \textbf{D - Incorrect:} Transfer = sharing with another party
			\end{itemize}
		\end{block}
	\end{frame}
	
	\section{Question 99: COSO ERM Components}
	
	\begin{frame}{Q99: COSO ERM Risk Components}
		\begin{block}{Question}
			Which are risk components of COSO ERM framework? (Choose three)
		\end{block}
		\begin{enumerate}[(A)]
			\item Risk response
			\item Internal environment
			\item Business continuity
			\item Control activities
		\end{enumerate}
	\end{frame}
	
	\begin{frame}{Q99: Explanation}
		\begin{block}{Correct Answer: A, B, D}
			COSO ERM components: internal environment, objective settings, event identification, risk assessment, risk response, control activities, information/communication, monitoring.
		\end{block}
		\begin{block}{Option Analysis}
			\begin{itemize}
				\item \textbf{A - Correct:} Risk response is component
				\item \textbf{B - Correct:} Internal environment is component
				\item \textbf{C - Incorrect:} Business continuity is NOT component
				\item \textbf{D - Correct:} Control activities is component
			\end{itemize}
		\end{block}
	\end{frame}
	
	\section{Question 100: Quantitative Analysis Output}
	
	\begin{frame}{Q100: Quantitative Risk Analysis Update}
		\begin{block}{Question}
			Which component is likely updated in risk register based on quantitative analysis?
		\end{block}
		\begin{enumerate}[(A)]
			\item Listing of risk responses
			\item Risk ranking matrix
			\item Listing of prioritized risks
			\item Qualitative analysis outcomes
		\end{enumerate}
	\end{frame}
	
	\begin{frame}{Q100: Explanation}
		\begin{block}{Correct Answer: C}
			Quantitative analysis outcome creates listing of prioritized risks updated in risk register.
		\end{block}
		\begin{block}{Option Analysis}
			\begin{itemize}
				\item \textbf{A - Incorrect:} Responses from Plan Risk Response
				\item \textbf{B - Incorrect:} Matrix from qualitative
				\item \textbf{C - Correct:} Prioritized risks from quantitative
				\item \textbf{D - Incorrect:} Qualitative outcomes separate
			\end{itemize}
		\end{block}
	\end{frame}
	
	\section{Final Review}
	
	\begin{frame}{CRISC Key Takeaways}
		\begin{block}{Essential Concepts}
			\begin{itemize}
				\item Risk identification is iterative
				\item Controls must be assessed for objective achievement
				\item KRIs provide early warning signals
				\item Risk responses: Avoid, Transfer, Mitigate, Accept (negative); Exploit, Enhance, Share (positive)
				\item Quantitative: SLE = AV × EF, ALE = SLE × ARO
			\end{itemize}
		\end{block}
	\end{frame}
	
	\begin{frame}{Final Advice}
		\begin{block}{Exam Preparation}
			\centering
			\vspace{0.5cm}
			\textbf{Remember:}
			\begin{itemize}
				\item Understand the \textit{why} behind each concept
				\item Know CIA triad (Confidentiality, Integrity, Availability)
				\item Understand maturity levels 0-5
				\item Know roles (data owner vs custodian)
				\item Focus on ISACA's framework approach
			\end{itemize}
			\vspace{1cm}
			\textbf{Good Luck on Your CRISC Exam!}
		\end{block}
	\end{frame}
	
\end{document}